\chapter{Reconstruction from modules and fusion}
\label{cyrec:chapter}

\section{Recovering the charge of a Heisenberg module}

A tensor product of representations can determine a Hopf algebra
even when the underlying fields have no coproduct at a fixed level.
The Heisenberg algebra gives an explicit example. Its oscillator
representation is fixed, while the charge varies. Removing the
oscillators recovers the charge space, and tensoring the charge spaces
defines fusion at the original level. The finite-dimensional charge
spaces also retain the logarithmic exchange of the fields.

Work over $\mathbb C$ and fix a nonzero current level
$\kappa_{\mathrm H}$. The modes satisfy
\begin{equation}\label{cyrec:eq:heisenberg}
 [J_m,J_n]=m\kappa_{\mathrm H}\delta_{m+n,0},
 \qquad m,n\in\mathbb Z.
\end{equation}
Let $\mathscr F_{\kappa_{\mathrm H}}$ consist of modules with
locally nilpotent positive modes. Require also $J_nv=0$ for
all sufficiently large $n$, depending on $v$.
No semisimplicity or local finiteness is imposed on $J_0$.
Morphisms commute with every mode.

Put $A=\mathbb C[q]$ and
$P=\mathbb C[x_1,x_2,\ldots]$, with finite polynomial states.
For an $A$-module $M$, write $Q_M$ for the action of $q$ and set
\begin{equation}\label{cyrec:eq:fock}
 F(M)=P\otimes M,
 \qquad J_{-n}=x_n,
 \qquad J_n=n\kappa_{\mathrm H}\partial_{x_n}\quad(n>0),
 \qquad J_0=Q_M.
\end{equation}
Conversely, the highest space of a module $V$ is
\[
 \Omega_+(V)=\bigcap_{n>0}\ker J_n,
\]
with charge $J_0$. The word highest refers only to the positive
oscillator modes. The charge can have arbitrary Jordan blocks.

\begin{theorem}\label{cyrec:thm:oscillators}
The functors $F$ and $\Omega_+$ are inverse exact equivalences
between $A\text{-}\mathrm{Mod}$ and
$\mathscr F_{\kappa_{\mathrm H}}$. The equivalence preserves all
fields obtained from polynomial Heisenberg states by derivatives and
normal products.
\end{theorem}

\begin{proof}
On $V$, write $X_n=J_{-n}$ and
$D_n=J_n/(n\kappa_{\mathrm H})$, so
$[D_n,X_m]=\delta_{nm}$. Define the projection
\begin{equation}\label{cyrec:eq:projector}
 E_n=\sum_{r\geq0}\frac{(-1)^r}{r!}X_n^rD_n^r,
 \qquad E=\prod_{n>0}E_n.
\end{equation}
Each sum is finite on a vector. Only finitely many factors of the
product act nontrivially, because different oscillator pairs commute
and the large positive modes annihilate that vector.
The Weyl relation gives $D_nE_n=0$, $E_nX_n=0$, and
$E_n|_{\ker D_n}=1$. Thus $E$ projects onto $\Omega_+(V)$.

For a finitely supported multi-index $\alpha$, put
$X^\alpha=\prod_nX_n^{\alpha_n}$,
$D^\alpha=\prod_nD_n^{\alpha_n}$ and
$\alpha!=\prod_n\alpha_n!$. The inverse to multiplication is
\begin{equation}\label{cyrec:eq:taylor}
 V\longrightarrow P\otimes\Omega_+(V),
 \qquad
 v\longmapsto\sum_\alpha\frac{x^\alpha}{\alpha!}\otimes ED^\alpha v.
\end{equation}
Commuting locally nilpotent derivatives make this sum finite.
The one-variable identity
\[
 \sum_{r\geq0}X^rE_nD_n^r/r!=1
\]
proves that one composite is the identity. For $w\in\Omega_+(V)$, the equality
$ED^\alpha X^\beta w=\delta_{\alpha\beta}\alpha!w$ proves the other.
These maps commute with $J_0$ as well as the oscillators.

The equivalence is additive, hence exact between the two abelian
categories. A mode-intertwining map intertwines derivatives and
normal products of the current. Such products define every
polynomial Heisenberg field, proving the last assertion.
\end{proof}

\subsection{The action monad and its tensor product}

Forget the zero mode for a moment. The same proof identifies the
oscillator category $\mathcal E_{\kappa_{\mathrm H}}$ with
$\mathrm{Vect}$. Its tensor is the balanced tensor over $P$:
\[
 (P\otimes M)\otimes_P(P\otimes N)=P\otimes(M\otimes N).
\]
Positive modes act by the sum of the two derivatives, which descends
to the balanced tensor, and negative modes multiply the common
polynomial factor. The unit is $P$. This construction uses a single
oscillator system at level $\kappa_{\mathrm H}$.

Adjoining an unconstrained zero mode is free. The left adjoint to
forgetting that mode sends $P\otimes M$ to
$P\otimes A\otimes M$. It defines the monad
\[
 \mathbb T(P\otimes M)=P\otimes A\otimes M,
\]
whose multiplication and unit come from multiplication and unit in
$A$. A module over this monad is exactly an action of $q$ on $M$.
The primitive Hopf structure
\begin{equation}\label{cyrec:eq:primitive-charge}
 \Delta q=q\otimes1+1\otimes q,
 \qquad\epsilon(q)=0,
 \qquad S(q)=-q
\end{equation}
makes $\mathbb T$ a comonoidal monad.

For a comonoidal monad, the left and right fusion maps combine its
comonoidal map with multiplication on the appropriate factor of
$\mathbb T^2$. It is called a \emph{Hopf monad} when both fusion
maps are invertible.
Here the two copies of $A$ can be written as $\mathbb C[u,v]$.
The fusion maps and their inverses are
\begin{equation}\label{cyrec:eq:monad-fusion}
\begin{array}{c|c|c}
 &\text{fusion map on }f(u,v)&\text{inverse on }f(u,v)\\ \hline
 \text{left}&f(u+v,v)&f(u-v,v)\\
 \text{right}&f(u+v,u)&f(v,u-v).
\end{array}
\end{equation}
Both assertions follow by substitution, over the polynomial ring
itself.

\begin{proposition}\label{cyrec:prop:monad}
The monad $\mathbb T$ is a Hopf monad. Its category of modules
has tensor and unit
\begin{equation}\label{cyrec:eq:charge-tensor}
 F(M)\star F(N)=F(M\otimes N),
 \qquad Q_{M\otimes N}=Q_M\otimes1+1\otimes Q_N,
 \qquad\mathbf1_\star=F(\mathbb C_0).
\end{equation}
When $M$ is finite-dimensional, the dual object is $F(M^*)$,
with charge $-Q_M^{\mathsf t}$.
\end{proposition}

\begin{proof}
Equation~\eqref{cyrec:eq:monad-fusion} verifies the two Hopf
conditions. The tensor of monad modules is induced by
\eqref{cyrec:eq:primitive-charge}, giving the displayed charge sum.
Evaluation and coevaluation for $M$ are charge-intertwining precisely
with the indicated dual action. Their usual triangular identities
remain identities after adjoining the common oscillator factor.
\end{proof}

The finite-charge category, where $\dim M<\infty$, is closed under
this tensor and these duals. The free-action monad does not preserve
it. In fact the forgetful functor from finite-dimensional
$A$-modules to finite-dimensional vector spaces has no left adjoint.
If its value $L$ on $\mathbb C$ existed, adjunction would give
$\operatorname{Hom}_A(L,\mathbb C_\lambda)=\mathbb C$ for every
$\lambda\in\mathbb C$. A finite-dimensional $L$ has finite charge
spectrum, so this Hom space vanishes for all other $\lambda$.
Finite-charge reconstruction therefore requires a different
algebraic construction.

\section{Matrix coefficients and the finite dual}

An action algebra describes operators on representations. Matrix
coefficients describe the representations as comodules. Their
universal quotient imposes every morphism in the category, so it
retains information that a chosen exchange matrix alone may omit.

Let $k$ be a field and $\mathcal C$ an essentially small
$k$-linear abelian category. Let
$\omega:\mathcal C\to\mathrm{Vect}_{\mathrm{fd}}(k)$ be faithful
and exact. Choose a small skeleton. Its matrix-coefficient space is
the quotient
\begin{equation}\label{cyrec:eq:coend}
 H_\omega=
 \left(\bigoplus_{X\in\mathcal C}\omega(X)^*\otimes\omega(X)\right)
 \big/\bigl\langle
 [\lambda\omega(f)\otimes v]_X-
 [\lambda\otimes\omega(f)v]_Y
 \bigr\rangle,
\end{equation}
where $f:X\to Y$, $v\in\omega(X)$ and
$\lambda\in\omega(Y)^*$. This is the coend
$\int^X\omega(X)^*\otimes\omega(X)$ in vector spaces.
Its elements are finite sums of matrix coefficients.

\begin{proposition}\label{cyrec:prop:coefficients}
For a basis $(e_i)$ of $\omega(X)$ and its dual $(e^i)$, the
formulas
\begin{align}
 \Delta[\lambda\otimes v]_X
 &=\sum_i[\lambda\otimes e_i]_X
                \otimes[e^i\otimes v]_X,\label{cyrec:eq:coend-delta}\\
 \epsilon[\lambda\otimes v]_X&=\lambda(v),
 \qquad
 \delta_X(v)=\sum_i e_i\otimes[e^i\otimes v]_X
 \label{cyrec:eq:coaction}
\end{align}
define a coalgebra and natural right coactions on $\omega(X)$.
\end{proposition}

\begin{proof}
The tensor $\sum_i e_i\otimes e^i$ is the identity endomorphism,
so the formulas are basis-independent. Inserting it twice gives
the same triple coefficient for either iterated coproduct.
Evaluation proves the counit and coaction identities.
For $f:X\to Y$, the relation in \eqref{cyrec:eq:coend} is exactly
the assertion that $\omega(f)$ intertwines the coactions. Inserting
the identity before or after $\omega(f)$ shows that the same
relation is preserved by $\Delta$ and $\epsilon$.
\end{proof}

\begin{theorem}[Reconstruction from an exact fiber functor]
\label{cyrec:thm:tannaka}
The canonical functor
\[
 \mathcal C\longrightarrow\operatorname{Comod}_{\mathrm{fd}}(H_\omega)
\]
is an equivalence. If $\mathcal C$ is rigid monoidal and $\omega$
is strong monoidal, $H_\omega$ is a Hopf algebra and the equivalence
is monoidal. If $\mathcal C$ is braided, its braiding determines a
coquasitriangular form on $H_\omega$ and the equivalence is braided.
\end{theorem}

\begin{proof}
The coalgebra reconstruction theorem for an exact faithful functor
to finite-dimensional vector spaces gives the first equivalence
\cite[Theorem~1.10.1]{EGNOTensorCategories}. Its universal coalgebra
is the matrix-coefficient quotient \eqref{cyrec:eq:coend}.
The hypotheses imply finite-dimensional Hom spaces and finite
length. Faithfulness embeds each Hom space into a finite matrix
space. Exactness sends every strict subobject chain to a strict
chain of subspaces. No separate comonadic hypothesis is
needed.

Write $J_{X,Y}:\omega(X)\otimes\omega(Y)\to\omega(X\otimes Y)$
for the tensor constraint. Multiplication is
\begin{equation}\label{cyrec:eq:coend-product}
 [\lambda\otimes v]_X[\mu\otimes w]_Y
 =[(\lambda\otimes\mu)J_{X,Y}^{-1}
                         \otimes J_{X,Y}(v\otimes w)]_{X\otimes Y}.
\end{equation}
Naturality of $J$ makes it well defined. The associativity and unit
constraints make it an associative unital product, and
\eqref{cyrec:eq:coend-delta} makes it a coalgebra map.
Dual evaluation and coevaluation give the antipode identities.
This is the rigid reconstruction of
\cite[Theorem~5.4.1]{EGNOTensorCategories}.

For a braiding $c_{X,Y}$, the form is specified by
\begin{equation}\label{cyrec:eq:coend-form}
 \rho([\lambda\otimes v]_X,[\mu\otimes w]_Y)
 = (\mu\otimes\lambda)J_{Y,X}^{-1}
                  \omega(c_{X,Y})J_{X,Y}(v\otimes w).
\end{equation}
Naturality descends it through the coefficient relations.
The two hexagons give its two multiplicativity identities, and
invertibility of $c$ gives its convolution inverse. On right
comodules it recovers
$v\otimes w\mapsto\sum w_{(0)}\otimes v_{(0)}
\rho(v_{(1)},w_{(1)})$, hence the original braiding.
\end{proof}

\subsection{The completion of natural endomorphisms}

The vector-space dual of a coend is an end. Accordingly,
\[
 H_\omega^*=\operatorname{End}(\omega)
\]
as algebras, with natural endomorphisms multiplied by composition.
For infinitely many matrix coefficients, this dual generally needs
a completed tensor product to carry its coproduct.

\begin{example}\label{cyrec:ex:torus}
Let $\mathcal C=\operatorname{Rep}_{\mathrm{alg}}(\mathbb G_m)$
and let $\omega$ forget the action. Every representation is a finite
sum of integral weight spaces. Naturality forces an endomorphism
of $\omega$ to be scalar on each weight, with independently chosen
scalars. Thus
\[
 \operatorname{End}(\omega)=\prod_{n\in\mathbb Z}k,
 \qquad (\Delta f)(m,n)=f(m+n).
\]

For $f=\delta_0$, restriction of this matrix to
$-N\leq m,n\leq N$ has rank $2N+1$.
An element of the ordinary tensor product of two function spaces
is a finite sum $\sum_{i=1}^r a_i(m)b_i(n)$, whose every finite
submatrix has rank at most $r$. Therefore $\Delta\delta_0$ is not
in that ordinary tensor product.

The product topology instead gives the continuous coproduct into
$k^{\mathbb Z\times\mathbb Z}$, with
$\epsilon(f)=f(0)$ and $S(f)(n)=f(-n)$.
The ordinary coefficient Hopf algebra is
\[
 H_\omega=k[t,t^{-1}],\qquad
 \Delta t=t\otimes t,\quad\epsilon(t)=1,\quad S(t)=t^{-1}.
\]
Its finite-dimensional comodules recover the original category.
\end{example}

Likewise, an exchange equation constrains a presentation without
necessarily presenting the coend. For the one-dimensional object
$X=k$ in $\mathrm{Vect}_{\mathrm{fd}}$, its self-exchange equation
alone imposes no relation on its single matrix coefficient $t$.
Adjoining a dual makes $t$ invertible and gives $k[t,t^{-1}]$.
The actual morphism $k=\mathbf1\to X$ imposes $t=1$ in
\eqref{cyrec:eq:coend}, and the full coend is $k$.
An RTT presentation therefore also needs the relations expressing
the morphisms and tensor constraints of the specified category.

\subsection{Finite charge and distributions}

Return to $A=\mathbb C[q]$. Its finite dual $A^\circ$ consists
of linear functionals annihilating an ideal of finite codimension.
Every such ideal is generated by a nonzero polynomial. Factoring
that polynomial and using the Chinese remainder theorem gives
\begin{equation}\label{cyrec:eq:finite-dual}
 A^\circ=\bigoplus_{\lambda\in\mathbb C}
               \bigoplus_{n\geq0}\mathbb C d_{\lambda,n},
 \qquad d_{\lambda,n}(p)=\frac{p^{(n)}(\lambda)}{n!}.
\end{equation}
Each element has finite support in both indices.

\begin{proposition}\label{cyrec:prop:finite-dual}
The finite dual is an ordinary Hopf algebra with
\begin{align}
 \Delta d_{\lambda,n}
  &=\sum_{i+j=n}d_{\lambda,i}\otimes d_{\lambda,j},
 &\epsilon(d_{\lambda,n})&=\delta_{n0},\label{cyrec:eq:distribution-delta}\\
 d_{\lambda,m}d_{\mu,n}
  &=\binom{m+n}{m}d_{\lambda+\mu,m+n},
 &S(d_{\lambda,n})&=(-1)^n d_{-\lambda,n}.
 \label{cyrec:eq:distribution-product}
\end{align}
Its unit is $d_{0,0}$. Its finite-dimensional comodules are the
finite-dimensional $A$-modules, and all its comodules are the
locally finite $A$-modules.
\end{proposition}

\begin{proof}
The Leibniz rule gives the coproduct. Differentiating
$p(u+v)$ at $(\lambda,\mu)$ gives the binomial product, and
substitution $q\mapsto-q$ gives the antipode.
These are the duals of the multiplication, coproduct and antipode
of $A$, and their values remain in the displayed direct sums.

On a generalized eigenspace $M_\lambda$, write
$Q_M=\lambda+N_\lambda$, with $N_\lambda$ nilpotent. The coaction
is
\begin{equation}\label{cyrec:eq:jordan-coaction}
 v\longmapsto\sum_{n\geq0}N_\lambda^n v\otimes d_{\lambda,n}.
\end{equation}
It is finite on $v$ and recovers the action by evaluation against
$p(q)$. Conversely, a coalgebra coaction places each vector in a
finite-dimensional subcomodule, and evaluation gives a locally
finite $A$-action. These two constructions are inverse and preserve
the primitive charge tensor.
\end{proof}

The same Hopf algebra is the coend for the forgetful functor on
finite-dimensional $A$-modules. Indeed their matrix coefficients
are exactly the functionals in \eqref{cyrec:eq:finite-dual}.
Forgetting an $A^\circ$-coaction has cofree right adjoint
$V\mapsto V\otimes A^\circ$. This is distinct from the free
left adjoint $V\mapsto A\otimes V$ on arbitrary $A$-modules.

\section{Logarithmic exchange and changes of presentation}

The primitive charge tensor has more than one possible braiding
on finite-dimensional modules. The endpoint operation determines
the relevant one. Give $x_n$ oscillator weight $n$ and take the
product completion in output oscillator weight. Inputs remain
polynomial. In the endpoint constructions, $M,N,L$ are
finite-dimensional charge modules. Write
\[
 J(z)=\sum_{m\in\mathbb Z}J_mz^{-m-1}=J^-(z)+J^+(z),
\]
where $J^-$ uses $m<0$ and $J^+$ uses $m\geq0$.

For a ground vector $v\in M$, define
\begin{equation}\label{cyrec:eq:endpoint}
 K_{M,N}(z)(v\otimes f(x)w)
 =\exp\left(\frac{Q_M}{\kappa_{\mathrm H}}
                       \sum_{r>0}\frac{x_rz^r}{r}\right)
       f(x_r-Q_Mz^{-r})(v\otimes w).
\end{equation}
In every output weight the exponential has finitely many terms.
Extend to the first input by
\begin{equation}\label{cyrec:eq:descendants}
 K(J_{-n}a,z)=
 \frac{\partial_z^{n-1}J^-(z)}{(n-1)!}K(a,z)
 +K(a,z)\frac{\partial_z^{n-1}J^+(z)}{(n-1)!}.
\end{equation}
The currents act on the output and on the second input,
respectively. Positive-mode sums are finite on an input, and
negative-mode sums are finite at each output weight.

\begin{proposition}\label{cyrec:prop:endpoint}
The operation $K$ satisfies the Ward identities
\begin{equation}\label{cyrec:eq:ward}
 J_mK(a,z)-K(a,z)J_m
 =\sum_{j\geq0}\binom mj z^{m-j}K(J_ja,z).
\end{equation}
Every holomorphic operation with these identities and
\eqref{cyrec:eq:descendants}, valued in the weight completion of
$F(L)$, is uniquely $F(\theta(z))K_{M,N}$.
Here $\theta(z):M\otimes N\to L$ is a holomorphic charge-intertwining
map on the chosen open subset of $\mathbb C^*$.
\end{proposition}

\begin{proof}
Commuting a positive mode through the exponential differentiates it
and gives $z^mQ_M$. Commuting a negative mode through the shift
gives $z^mQ_M$ as well. The zero-mode equation is charge additivity.
These are \eqref{cyrec:eq:ward} for a highest first input.
Normal ordering gives the descendant equation and extends the Ward
identity by induction on its oscillator degree.

Conversely, remove the exponential and shift from a highest-input
operation. The remaining map commutes with every oscillator
multiplication and derivative, so it is the extension of its ground
map. The zero-mode identity makes that map charge-intertwining.
The descendant equation then determines all remaining values.
\end{proof}

The translation on $F(M)$ is
\begin{equation}\label{cyrec:eq:translation}
 T_M=\sum_{n>0}n x_{n+1}\partial_{x_n}
                       +\frac{x_1Q_M}{\kappa_{\mathrm H}}.
\end{equation}
It satisfies $[T_M,J_m]=-mJ_{m-1}$, including $m=1$, where the
last term supplies $-Q_M$. Put
$t_{M,N}=Q_M\otimes Q_N/\kappa_{\mathrm H}$.
Direct differentiation of \eqref{cyrec:eq:endpoint} gives
\begin{equation}\label{cyrec:eq:translation-defect}
 T_{M\otimes N}K-KT_N-\partial_zK=z^{-1}t_{M,N}K.
\end{equation}
The oscillator terms cancel consecutively, leaving
$Q_MQ_N/(\kappa_{\mathrm H}z)$.

Thus the connection on representing states and the connection on
operations have opposite signs:
\begin{equation}\label{cyrec:eq:connections}
 \nabla^{\mathrm{state}}=d+t_{M,N}\,d\log z,
 \qquad
 \nabla^{\mathrm{op}}\theta=d\theta-\theta t_{M,N}\,d\log z.
\end{equation}
Their horizontal factors are $z^{-t_{M,N}}$ and $z^{t_{M,N}}$.
For finite charge spaces these are ordinary matrix exponentials
after a branch of $\log z$ is chosen.
Denote the normalized horizontal endpoint operation by
$\mathcal E_{M,N}=F(z^{t_{M,N}})K_{M,N}$.

\begin{proposition}\label{cyrec:prop:exchange}
The normalized horizontal endpoint operations have braiding
\begin{equation}\label{cyrec:eq:braiding}
 c^+_{M,N}=\tau\exp\left(\frac{\pi i}{\kappa_{\mathrm H}}
                                      Q_M\otimes Q_N\right),
\end{equation}
where the half-turn of the separation is counterclockwise and
$\tau$ exchanges the two charge factors. Their associator is the
usual identification of the threefold charge tensor.
\end{proposition}

\begin{proof}
For three charges put $t_{ij}=Q_iQ_j/\kappa_{\mathrm H}$.
For inputs $u_i\in F(M_i)$, the ordered and nested operations are
\begin{align*}
 \Phi_{23}&=\mathcal E_{M_1,M_2\otimes M_3}(u_1,z_1)
                          \mathcal E_{M_2,M_3}(u_2,z_2)u_3,\\
 \Phi_{12}&=\mathcal E_{M_1\otimes M_2,M_3}
             (\mathcal E_{M_1,M_2}(u_1,z_1-z_2)u_2,z_2)u_3.
\end{align*}
They are expanded in $|z_1|>|z_2|>0$ and
$|z_2|>|z_1-z_2|>0$, respectively.

All pair residues commute. Reordering the creation and
annihilation factors in the ordered channel gives
\[
 z_1^{t_{12}+t_{13}}(1-z_2/z_1)^{t_{12}}z_2^{t_{23}}.
\]
The nested channel gives
\[
 (1+(z_1-z_2)/z_2)^{t_{13}}
      z_2^{t_{13}+t_{23}}(z_1-z_2)^{t_{12}}.
\]

With consistent logarithms these are both
$z_1^{t_{13}}z_2^{t_{23}}(z_1-z_2)^{t_{12}}$.
Their series converge on $|z_1|>|z_2|>|z_1-z_2|>0$.
The Ward and descendant recursions reduce every further coefficient
to this ground expression, by finite operations in each output weight.
Their rational coefficients have poles only at the three collisions. Hence the equality extends to every polynomial input.

The exchanged channel replaces $z_1-z_2$ by $z_2-z_1$.
A counterclockwise half-turn contributes $\exp(\pi i t_{12})$.
Charge additivity and the commuting residues prove both hexagons
for \eqref{cyrec:eq:braiding}. The zero charge proves its two unit
identities. This compares the separated and nested operations on
their logarithmic domains.
\end{proof}

Under Proposition~\ref{cyrec:prop:finite-dual}, this braiding is
represented by the scalar-valued form
\begin{equation}\label{cyrec:eq:distribution-form}
 \rho^+(d_{\lambda,m},d_{\mu,n})
 =\frac{1}{m!n!}\partial_\lambda^m\partial_\mu^n
                  \exp(\pi i\lambda\mu/\kappa_{\mathrm H}).
\end{equation}
Taylor expansion on each pair of generalized eigenspaces gives
\eqref{cyrec:eq:braiding}. The two exponential addition identities
give the coquasitriangular identities, and changing the sign of the
exponent gives the convolution inverse. The scalar exponential
therefore presents the full finite-charge category, including
nonsemisimple charges.

\begin{proposition}\label{cyrec:prop:ambient-braiding}
On all $\mathbb C[q]$-modules with the primitive charge tensor, the
only natural braiding satisfying the unit identities is the flip.
\end{proposition}

\begin{proof}
Remove the flip from a proposed braiding. On the pair of free
modules $(A,A)$ its value on $1\otimes1$ is a polynomial
$r(u,v)$. Naturality with respect to the module maps
$A\to M$, $1\mapsto m$, shows that the transformation on every
pair is the action of this polynomial. The inverse transformation
likewise comes from a polynomial, so $r$ is a unit in
$\mathbb C[u,v]$. It is constant. The unit identities set that
constant equal to one.
\end{proof}

\subsection{The full turn and a nonsplit extension}

A monoidal change of presentation with tensorator $J_{M,N}$
changes the braiding and its full turn by
\begin{equation}\label{cyrec:eq:conjugacy}
 c'_{M,N}=J_{N,M}^{-1}c_{M,N}J_{M,N},
 \qquad
 D'_{M,N}=J_{M,N}^{-1}D_{M,N}J_{M,N},
\end{equation}
where
\[
 D_{M,N}=c^+_{N,M}c^+_{M,N}
        =\exp(2\pi iQ_M\otimes Q_N/\kappa_{\mathrm H}).
\]
Consequently a nonidentity full turn cannot become the identity
through an equivalence preserving this exchange.

For a scalar example take $\kappa_{\mathrm H}=1$ and charges
$\lambda=\mu=1/\sqrt2$. Then $D=-1$.

There is also an example with only zero eigenvalues. On
$M=\mathbb C e_0\oplus\mathbb C e_1$, let
$Ne_1=e_0$ and $Ne_0=0$. The charge module with $Q_M=N$ is the
nonsplit self-extension
\[
 0\longrightarrow\mathbb C_0\longrightarrow M
                     \longrightarrow\mathbb C_0\longrightarrow0.
\]
It cannot split because a charge-zero complement would lie in
$\ker N=\mathbb C e_0$.

On its tensor square,
\begin{equation}\label{cyrec:eq:nilpotent-turn}
 D_{M,M}=1+\frac{2\pi i}{\kappa_{\mathrm H}}N\otimes N\ne1,
 \qquad
 z^{N\otimes N/\kappa_{\mathrm H}}
 =1+\frac{\log z}{\kappa_{\mathrm H}}N\otimes N.
\end{equation}
The nonzero rank-one operator $N\otimes N$ detects both the
logarithm and its monodromy. Replacing a charge by its eigenvalues
would lose this extension and its exchange.

\begin{proposition}\label{cyrec:prop:symmetric-sectors}
Let $\mathcal D$ be a full replete subcategory of finite-dimensional
charge modules, containing the unit and closed under finite sums,
subquotients, tensor products and duals. The braiding \eqref{cyrec:eq:braiding} is symmetric on $\mathcal D$
exactly under the following conditions. Every charge is semisimple,
and the simple charges form an additive subgroup $\Gamma\subset\mathbb C$
satisfying
\begin{equation}\label{cyrec:eq:integral-charges}
 \lambda\mu/\kappa_{\mathrm H}\in\mathbb Z
                \qquad(\lambda,\mu\in\Gamma).
\end{equation}
Every nonzero such subgroup is $\Gamma=\gamma\mathbb Z$, with
$\gamma^2/\kappa_{\mathrm H}\in\mathbb Z\setminus\{0\}$.
\end{proposition}

\begin{proof}
Suppose a charge has a nontrivial Jordan block at $\lambda$.
A two-dimensional subquotient has charge $\lambda+N$, with
$N^2=0$ and $N\ne0$. If $\lambda\ne0$, its full turn with the
simple subquotient $\mathbb C_\lambda$ contains the nonzero
nilpotent term $2\pi i\lambda N/\kappa_{\mathrm H}$ after its
scalar factor. It cannot be the identity. If $\lambda=0$,
\eqref{cyrec:eq:nilpotent-turn} gives the contradiction.
Thus all charges are semisimple.

Closure under tensor and dual makes the simple charges an additive
subgroup. On two simples, symmetry is exactly
\eqref{cyrec:eq:integral-charges}. Conversely, that condition makes
the full turn the identity on every sum of these simples.

Choose $\lambda_0\ne0$ in $\Gamma$. The injective map
$\lambda\mapsto\lambda\lambda_0/\kappa_{\mathrm H}$ embeds
$\Gamma$ as a subgroup of $\mathbb Z$. Its image is $d\mathbb Z$
for some positive integer $d$, so $\Gamma=\gamma\mathbb Z$.
Applying \eqref{cyrec:eq:integral-charges} to $\gamma$ proves the
last assertion.
\end{proof}

These semisimple subcategories need not be closed under extensions
in the ambient charge category. The preceding neutral extension
already prevents that closure. For $\Gamma=\gamma\mathbb Z$,
the coefficient Hopf algebra is $\mathbb C[t,t^{-1}]$ and its
form is $\rho(t^m,t^n)=(-1)^{Nmn}$, where
$N=\gamma^2/\kappa_{\mathrm H}$. This describes a symmetric
charge subcategory. It does not identify an arbitrary category of
fields with a spectral presentation.

\section{The affine resolution of a dg line algebra}

The charge equivalence separates a fixed oscillator representation
from its multiplicity space. For a nonabelian current algebra the
corresponding construction uses a resolution, and its endomorphisms
form a dg algebra. The opposite algebra and the shifted affine level
are both forced by the action on that resolution.

Use the polynomial dg algebra $\mathcal Y_k$ of
Definition~\ref{cylev:def:line-algebra}, for a finite-dimensional
complex Lie algebra $\mathfrak g$ with invariant symmetric form
$\kappa$. The scalar $k$ is its level parameter. Put
\[
 h_{\mathfrak g}(x,y)=\tfrac12\operatorname{tr}
                  (\operatorname{ad}x\,\operatorname{ad}y),
 \qquad \ell=-k\kappa-h_{\mathfrak g}.
\]

Let $A_\ell$ be the affine enveloping algebra with central element
acting as one, so that
\begin{equation}\label{cyrec:eq:affine}
 [xt^r,yt^s]=[x,y]t^{r+s}
                      +r\delta_{r+s,0}\ell(x,y)1.
\end{equation}
Write $\mathfrak n=t\mathfrak g[t]$ and
$\mathfrak l=\mathfrak g[t^{-1}]$. The induced module
\[
 \mathbb W_\ell=A_\ell\otimes_{U(\mathfrak n)}\mathbb C
\]
has underlying vector space $U(\mathfrak l)$ by PBW. Its cyclic
vector is killed by $\mathfrak n$, while its horizontal zero modes
act freely.

\subsection{A resolution with a right action}

The induced Chevalley resolution has
\begin{equation}\label{cyrec:eq:affine-resolution}
 \mathcal P^{-p}=A_\ell\otimes\Lambda^p\mathfrak n,
 \qquad p\geq0,
\end{equation}
with augmentation $\pi:\mathcal P\to\mathbb W_\ell$ and
differential
\begin{align}
 d_{\mathcal P}(u\otimes y_1\wedge\cdots\wedge y_p)
 ={}&\sum_i(-1)^{i-1}uy_i\otimes
              y_1\wedge\cdots\widehat y_i\cdots\wedge y_p
 \nonumber\\
 &+\sum_{i<j}(-1)^{i+j}u\otimes[y_i,y_j]\wedge
              y_1\wedge\cdots\widehat y_i\cdots\widehat y_j
                       \cdots\wedge y_p.
 \label{cyrec:eq:chevalley-differential}
\end{align}
The surviving exterior factors retain their order. PBW makes
$A_\ell$ free as a right $U(\mathfrak n)$-module, so inducing the
free Chevalley resolution proves that $\pi$ is a quasi-isomorphism.
Every chain is an algebraic finite sum.

Give $xt^r$ energy $-r$, including when it is an exterior generator.
The differential preserves energy, and $\mathbb W_\ell$ has
nonnegative energy. Define
\begin{equation}\label{cyrec:eq:endomorphisms}
 \mathcal E=\bigoplus_{w\in\mathbb Z}
      \operatorname{End}_{A_\ell}^{\bullet,w}(\mathcal P),
 \qquad
 d_{\mathcal E}(F)=d_{\mathcal P}F-(-1)^{|F|}Fd_{\mathcal P}.
\end{equation}
Here $w$ is the internal energy shift. Each Hom complex uses the
usual product totalization over the cohomological terms of
$\mathcal P$, while the outer sum allows only finitely many energy
shifts. The graded opposite product is
$F\cdot_{\mathrm{op}}G=(-1)^{|F||G|}GF$.

For $x\in\mathfrak l$, let $R_x$ multiply the affine factor on
the right, and let $a_x$ be the exterior derivation induced by
$y\mapsto\operatorname{pr}_{\mathfrak n}[x,y]$.
Projection is along $\mathfrak l$ and the central direction.
For the basis $(t_a)$ of $\mathfrak g$, let $\iota_{a,s}$
contract with the functional dual to $t_at^s$, for $s\geq1$.
Set
\begin{equation}\label{cyrec:eq:affine-lifts}
 \beta_{a,n}=R_{t_at^{-n}}-a_{t_at^{-n}},
 \qquad\gamma_{a,m}=\iota_{a,m+1},
 \qquad n,m\geq0.
\end{equation}
Their cohomological degrees are zero and one, and their energies
are $n$ and $m+1$.

\begin{lemma}\label{cyrec:lem:affine-lifts}
With ordinary operator composition, the lifts satisfy
\begin{align}
 [\beta_{a,n},\beta_{b,m}]&=-f_{ab}^{\ c}\beta_{c,n+m},
 & [\beta_{a,n},\gamma_{b,m}]&=f_{ac}^{\ b}\gamma_{c,n+m},
 \label{cyrec:eq:lift-brackets}\\
 [d_{\mathcal P},\gamma_{a,n}]
 &=\tfrac12f_{bc}^{\ a}\sum_{m+p=n-1}\gamma_{b,m}\gamma_{c,p},
 \label{cyrec:eq:lift-ghost}\\
 [d_{\mathcal P},\beta_{a,n}]
 &=-n\ell_{ab}\gamma_{b,n-1}
       +f_{ab}^{\ c}\sum_{m+p=n-1}\beta_{c,p}\gamma_{b,m}.
 \label{cyrec:eq:lift-even}
\end{align}
The contractions anticommute, and negative indices are absent.
\end{lemma}

\begin{proof}
Right multiplications obey $[R_x,R_y]=-R_{[x,y]}$.
The projected exterior action is a Lie action and commutes with
right multiplication. These facts give the first bracket.
The identity $[a_x,\iota_\lambda]=-\iota_{\lambda\circ a_x}$
gives the second.

Writing $\varepsilon_y$ for exterior multiplication, the two
pieces of the differential are
\[
 d_1=\sum_iR_{y_i}\iota_i,
 \qquad
 d_2=\tfrac12\sum_{i,j}\varepsilon_{[y_i,y_j]}\iota_i\iota_j.
\]
Only contractions from the input exterior monomial contribute.
The anticommutator with $\gamma$ comes from $d_2$ and gives
\eqref{cyrec:eq:lift-ghost}. In $[d_1,\beta_{a,n}]$ the
positive-mode terms cancel the projected exterior action.
The remaining current and central terms, for $1\leq s\leq n$,
are $f_{ab}^{\ c}t_ct^{s-n}-n\delta_{s,n}\ell_{ab}$.

Jacobi, projected to $\mathfrak n$, gives
\[
 [d_2,\beta_x]=
       -\sum_i a_{\operatorname{pr}_{\mathfrak l}[x,y_i]}\iota_i.
\]
It supplies the exterior correction to those remaining currents.
Putting $s=m+1$ proves \eqref{cyrec:eq:lift-even}.
\end{proof}

\begin{theorem}\label{cyrec:thm:affine-endomorphisms}
The map
\begin{equation}\label{cyrec:eq:rho}
 \rho:\mathcal Y_k\longrightarrow\mathcal E^{\mathrm{op}},
 \qquad B^a_n\longmapsto\beta_{a,n},
 \qquad c_{a,n}\longmapsto\gamma_{a,n}
\end{equation}
is a unital energy-preserving quasi-isomorphism of dg algebras.
In particular,
\[
 \mathcal Y_k\simeq
 R\operatorname{Hom}^{\mathrm{wt}}_{A_\ell}
       (\mathbb W_\ell,\mathbb W_\ell)^{\mathrm{op}},
 \qquad\ell=-k\kappa-h_{\mathfrak g},
\]
where the superscript denotes the direct sum of homogeneous energy
shifts. No nondegeneracy of $\kappa$ is required.
\end{theorem}

\begin{proof}
The graded opposite changes the two brackets of
\eqref{cyrec:eq:lift-brackets} into the line relations and leaves
the product of two odd contractions unchanged.
For the even differential, invariance gives
$-\tfrac12f_{ab}^{\ c}f_{cd}^{\ b}=(h_{\mathfrak g})_{ad}$.
There are $n$ pairs with $m+p=n-1$, so symmetric reordering gives
\[
 [d_{\mathcal P},\beta_{a,n}]
 =-n(\ell+h_{\mathfrak g})_{ab}\gamma_{b,n-1}
   +\tfrac12 f_{ab}^{\ c}\sum_{m+p=n-1}
       (\beta_{c,p}\gamma_{b,m}+\gamma_{b,m}\beta_{c,p}).
\]
Since $-(\ell+h_{\mathfrak g})=k\kappa$, this is the image of
the line differential. Thus $\rho$ is a dg algebra map.

Compose $\rho$ with augmentation of its output:
\begin{equation}\label{cyrec:eq:cochain-evaluation}
 \mathcal Y_k\longrightarrow\mathcal E^{\mathrm{op}}
 \xrightarrow{\pi\circ-}
 \bigoplus_w\operatorname{Hom}^{\bullet,w}_{A_\ell}
                         (\mathcal P,\mathbb W_\ell).
\end{equation}
The last target is used only as a complex.
Order the basis $(y_i)$ of $\mathfrak n$ and write its dual
ghosts $c_i$. On the ordered wedge
$y_{i_1}\wedge\cdots\wedge y_{i_p}$, the ghost-left PBW
monomial $c_{i_1}\cdots c_{i_p}u$, with
$u\in U(\mathfrak l)$, evaluates to
$(-1)^{p(p-1)/2}u|0\rangle$ and vanishes on other basis wedges
of length $p$.
The reversed composition applies the contractions first, after
which all exterior corrections vanish.

Fixing the energy shift $w$ makes this a basis of the entire
Hom complex. A wedge whose positive mode indices sum to $s$ has
energy $-s$, and its output must have energy $w-s\geq0$.
There are only finitely many such wedges. Their outputs remain
algebraic elements of $U(\mathfrak l)$, including its possibly
infinite-dimensional energy-zero part.
Hence the composite in \eqref{cyrec:eq:cochain-evaluation} is
an isomorphism of complexes.

The bounded-above graded-free resolution $\mathcal P$ is
$K$-projective in internally graded modules. Applying each
homogeneous Hom complex to $\pi$ makes the second arrow a
quasi-isomorphism. Direct sums over shifts preserve it, so $\rho$
is a quasi-isomorphism as well.
\end{proof}

For $\mathfrak{sl}_2$, with $\kappa(E,F)=1$ and
$\kappa(H,H)=2$, one has $h_{\mathfrak g}=2\kappa$ and
affine level $-k-2$. On the single exterior generator $Et$,
$[d_{\mathcal P},\beta_{F,1}]$ gives $-H+k+2$.
Using affine level $-k$ would instead lose the constant two and
fail to intertwine the unchanged line differential.

The quasi-isomorphism also has a useful limit on adjoining a formal
differential generator. Let
\[
 \mathcal Y_k^q=\mathcal Y_k\langle q\rangle/
     (q^2,\ qa-(-1)^{|a|}aq-Qa),
 \qquad |q|=1,
\]
with differential $[q,-]$ and energy zero for $q$.
The map extends by $q\mapsto-d_{\mathcal P}$.
However, if $N$ multiplies $\mathcal P^i$ by $i$, then
$d_{\mathcal E}N=-d_{\mathcal P}$. The charge image is exact,
whereas the energy-zero summand $U(\mathfrak g)q$ has zero
differential and nonzero class $[q]$ in the source.
The extended map is therefore not a quasi-isomorphism.

\subsection{The derived functor and the spectral comparison}

The opposite algebra in \eqref{cyrec:eq:rho} defines the right
action
\[
 p\cdot a=(-1)^{|p||a|}\rho(a)p
\]
on the left affine module $\mathcal P$. It gives a derived functor
\begin{equation}\label{cyrec:eq:derived-functor}
 \Phi_k(M)=\mathcal P\otimes^{\mathbb L}_{\mathcal Y_k}M
       \quad\text{in }D(A_\ell\text{-}\mathrm{Mod}).
\end{equation}
It sends the regular line module to $\mathbb W_\ell$.
Suppose the internally graded complex underlying $M$ has energy
bounded below. Every cohomology module of $\Phi_k(M)$ retains
that lower bound and is smooth for the affine action.

To prove the energy assertion, use the bar resolution over
$\mathbb C$. The line algebra has nonnegative energy, so each
free coefficient complex retains the lower bound of $M$.
After tensoring with $\mathcal P$, its free associated piece
can be replaced by $\mathbb W_\ell$, also of nonnegative
energy. Cones and the homotopy colimit of the finite bar stages
preserve vanishing below the bound. A sufficiently positive mode
sends any fixed-energy cohomology class below that bound, proving
smoothness.

When $\mathfrak g$ is abelian and $k\kappa$ is nondegenerate,
the positive-energy pairs in $\mathcal Y_k$ contract, so
$U(\mathfrak g)\hookrightarrow\mathcal Y_k$ is a quasi-isomorphism.
The zero-mode exterior correction vanishes and $\pi$ respects
the right $U(\mathfrak g)$-action. Consequently
\[
 \Phi_k(M)\simeq\mathbb W_{-k\kappa}
                         \otimes_{U(\mathfrak g)}M|_{U(\mathfrak g)}.
\]
PBW makes the right module free, so this ordinary tensor computes
the derived tensor. In rank one it is the Fock induction already
constructed above.

For the spectral comparison take $k\ne0$ and put
$K_z=\mathbb C((z^{-1}))$. The notation $V((z^{-1}))$ means
Laurent series bounded above in powers of $z$, with algebraic
coefficients in $V$. It is generally larger than $V\otimes K_z$.

Define the affine mode map
\begin{equation}\label{cyrec:eq:affine-coproduct}
 \delta_z(xt^m)=
 \sum_{j\geq0}\binom{-m}{j}z^{-m-j}xt^{-j}\otimes1
                       +1\otimes xt^m.
\end{equation}
Its target is $(U(\mathfrak l)\otimes A_\ell)((z^{-1}))$.
The first component expands $x(z+t^{-1})^{-m}$.
It has no affine cocycle, because the residue of $f\,dg$
vanishes for series in $t^{-1}$ with nonnegative powers of
$t^{-1}$. The second component retains the central term.
Thus \eqref{cyrec:eq:affine-coproduct} extends to an algebra map.

On algebraic mode inputs it satisfies
\begin{equation}\label{cyrec:eq:affine-coassociativity}
 (\delta_w\otimes1)\delta_z
       =(1\otimes\delta_z)\delta_{z+w},
 \qquad (\epsilon\otimes1)\delta_z=1.
\end{equation}
Both iterates send $xt^m$ to the sum of
$x(z+w+t^{-1})^{-m}$, $x(z+t^{-1})^{-m}$ and $xt^m$ in
the three respective factors. The counit here is the augmentation
of $U(\mathfrak l)$ in the first factor.

Use the line coproduct $\Delta_z$ and degree-one twist $r(z)$
of \eqref{cylev:eq:line-spectral}. On
$(\mathcal P\otimes\mathcal P)((z^{-1}))$ put
\begin{equation}\label{cyrec:eq:resolution-twist}
 D(v)=(d_{\mathcal P}\otimes1+1\otimes d_{\mathcal P})v
                              -(-1)^{|v|}v\cdot r(z).
\end{equation}
The right tensor action is
$(p_1\otimes p_2)(a_1\otimes a_2)
=(-1)^{|p_2||a_1|}p_1a_1\otimes p_2a_2$.
The Maurer--Cartan identity for $r$ gives $D^2=0$.

\begin{theorem}\label{cyrec:thm:resolution-coproduct}
The map
\begin{equation}\label{cyrec:eq:resolution-coproduct}
 \Delta_{\mathcal P}(u\otimes\omega)
       =\delta_z(u)\bigl(1\otimes(1\otimes\omega)\bigr)
\end{equation}
is a chain map from $\mathcal P$ to the twisted complex
\eqref{cyrec:eq:resolution-twist}. It intertwines the left
affine action along $\delta_z$ and the right line action along
$\Delta_z$. It is counital on the first factor and coassociative
on the common coefficient domain of the iterated spectral maps.
\end{theorem}

\begin{proof}
All exterior factors are placed in the second resolution.
The discrepancy between its differential and the image of
\eqref{cyrec:eq:chevalley-differential} consists of the first
components of $\delta_z(y_i)$, one for each exterior contraction.
The terms $B\otimes c$ in $r$ give precisely these components
with the sign in \eqref{cyrec:eq:resolution-twist}. The terms
$c\otimes B$ contract the first vacuum and vanish. This proves
$D\Delta_{\mathcal P}=\Delta_{\mathcal P}d_{\mathcal P}$.

For a right ghost the first translated contraction likewise
annihilates the first vacuum, leaving the contraction of $\omega$
in the second factor. For a right even mode, the two translated
right multiplications are the two components of
$\delta_z(t_at^{-n})$, and its exterior correction remains in
the second factor. These identities prove right linearity on
generators, hence on the entire polynomial line algebra.
Left linearity is built into the definition.

Both triple maps put the exterior factors in the last copy.
Equation~\eqref{cyrec:eq:affine-coassociativity} identifies their
affine factors, and Theorem~\ref{cylev:thm:line-spectral}
identifies their differential twists.
For shifted kernels this equality is read first on
$0<|w|<|z|<|z+w|$ with the stated small auxiliary jet variables,
then on their rationally continued mode coefficients.
It asserts no extension of translation to arbitrary iterated
Laurent completions. The first-factor augmentation proves the
counit identity.
\end{proof}

For left dg line modules $M,N$, let $M\otimes_zN$ be
$(M\otimes N)((z^{-1}))$ with differential $d+\rho_{M,N}(r)$
and line action through $\Delta_z$. Let
$X_1\circledast_zX_2=(X_1\otimes X_2)((z^{-1}))$ with affine
action through $\delta_z$. These specify two different
completed tensor constructions.

\begin{corollary}\label{cyrec:cor:tensorator}
There is a natural chain map
\begin{equation}\label{cyrec:eq:tensorator}
 \theta_{M,N}:\mathcal P\otimes_{\mathcal Y_k}(M\otimes_zN)
 \longrightarrow
 (\mathcal P\otimes_{\mathcal Y_k}M)
          \circledast_z(\mathcal P\otimes_{\mathcal Y_k}N),
 \qquad p\otimes v\longmapsto\Delta_{\mathcal P}(p)\otimes v.
\end{equation}
It is affine-linear and inherits the counit and coassociativity
identities on the coefficient domains of
Theorem~\ref{cyrec:thm:resolution-coproduct}.
Suppose $M,N$ are homotopically flat over $\mathcal Y_k$, and
$M\otimes_zN$ is homotopically flat over $\mathcal Y_k\otimes K_z$.
Then the map computes the comparison of derived images, with
algebraic scalar extension to $K_z$ on the source.
\end{corollary}

\begin{proof}
Right linearity of $\Delta_{\mathcal P}$ makes the map balanced.
The right and left twists cancel under this balancing, leaving
the ordinary tensor differential on the target. Multiplication
of Laurent coefficients is defined because each series has an
upper bound, so each output coefficient is a finite sum.
The remaining identities follow from the preceding theorem.
Under the flatness assumptions, the indicated tensors compute
derived tensors. Laurent completion of complexes of vector spaces
is exact, being a filtered union of products, and therefore
preserves the target quasi-isomorphisms.
\end{proof}

This target carries an affine action, but its Laurent completion
need not be smooth as an algebraic affine module. Nor does the
construction identify $\circledast_z$ with a separately defined
affine fusion product. Those are additional questions about the
category and the comparison map.

\subsection{Translation and its homotopy}

The coordinate change $t^{-1}\mapsto t^{-1}+z$ is generated by
$L_1=-t^2\partial_t$. Its action on the resolution is
\[
 T_{\mathcal P}=\operatorname{ad}L_1\otimes1+1\otimes L_1.
\]
It commutes with $d_{\mathcal P}$ because $L_1$ preserves the
bracket and its central cocycle. It is a derivation along the
affine action and the line translation, and
\begin{equation}\label{cyrec:eq:resolution-horizontal}
 \partial_z\Delta_{\mathcal P}
                     =(T_{\mathcal P}\otimes1)\Delta_{\mathcal P}.
\end{equation}
Thus $\partial_z-T_{\mathcal P}\otimes1$ commutes with $D$.

Assume now that $\kappa$ is nondegenerate. An infinite
Sugawara sum requires a completion. Put
$\mathfrak n_N=t^N\mathfrak g[t]$ and let $F_N\mathcal P$
be spanned by chains whose affine factor lies in
$A_\ell\mathfrak n_N$ or whose exterior factor contains
$\mathfrak n_N$. Define
\[
 \widehat{\mathcal P}=\varprojlim_N\mathcal P/F_N\mathcal P
\]
degreewise.

PBW identifies each quotient with
$U(\mathfrak l)$ tensored with the Chevalley resolution for
$\mathfrak n/\mathfrak n_N$. Its cohomology is
$\mathbb W_\ell$ in degree zero, and the transition maps induce
the identity there. The chain towers are surjective, so the
inverse-limit exact sequence proves
$\mathcal P\simeq\widehat{\mathcal P}$.

Right multiplication by $xt^{-m}$ and its exterior correction
send $F_N$ into $F_{N-m}$ for $N>m$. Fixed contractions are
continuous as well. The bimodule action therefore extends.
On this completion define
\begin{align}
 L_1^{\mathrm{Sug}}
 &=-\frac1k\sum_{m\geq0}\kappa^{ab}
                          t_at^{-m}t_bt^{m+1},
 \label{cyrec:eq:sugawara}\\
 L_{\mathcal P}
 &=\frac1k\sum_{m\geq0}\kappa^{ab}
    \left(R_{t_bt^{-m}}-\tfrac12a_{t_bt^{-m}}\right)
                          \varepsilon_{t_at^{m+1}}.
 \label{cyrec:eq:translation-homotopy}
\end{align}
The first acts by left multiplication and the second has
cohomological degree $-1$. Their inverse level form is
$(\ell+h_{\mathfrak g})^{-1}=-k^{-1}\kappa^{-1}$.

\begin{proposition}\label{cyrec:prop:translation-homotopy}
The sums converge continuously on $\widehat{\mathcal P}$ and
satisfy
\begin{align}
 [L_{\mathcal P},\beta_{a,n}]&=0,
 & [L_{\mathcal P},\gamma_{a,n}]_+
      &=k^{-1}\kappa^{ab}\beta_{b,n},
 \label{cyrec:eq:homotopy-actions}\\
 d_{\mathcal P}L_{\mathcal P}+L_{\mathcal P}d_{\mathcal P}
      &=T_{\mathcal P}-L_1^{\mathrm{Sug}}.
 \label{cyrec:eq:homotopy-identity}
\end{align}
\end{proposition}

\begin{proof}
The term of index $m$ in $L_{\mathcal P}$ lies in $F_{m+1}$.
The only exterior correction that could remove its new high
factor is proportional to $\kappa^{ab}[t_b,t_a]t=0$.
Modulo $F_d$ only $m<d$ remains, and those finitely many
operators send $F_{2d}$ into $F_d$. This proves continuity.
The same estimate after commuting a high positive mode past a
fixed algebraic input proves convergence of
$L_1^{\mathrm{Sug}}$.

Invariance of $\kappa^{-1}$ cancels the commutator with each
$\beta$. For a contraction, use
$[\varepsilon_y,\iota_\lambda]_+=\lambda(y)$ and
$[a_x,\iota_\lambda]=-\iota_{\lambda\circ a_x}$.
Contracting the new exterior factor gives
$k^{-1}\kappa^{ab}(R_{t_bt^{-n}}-\tfrac12a_{t_bt^{-n}})$.
The remaining contraction through the exterior correction gives
the other $-\tfrac12k^{-1}\kappa^{ab}a_{t_bt^{-n}}$.
This proves \eqref{cyrec:eq:homotopy-actions}.

Let $Z$ be the left side of \eqref{cyrec:eq:homotopy-identity}
minus its right side. The derivation identities make $Z$
left affine-linear and right line-linear, and it commutes with
the differential. On the cyclic generator,
$d_{\mathcal P}L_{\mathcal P}(1)=-L_1^{\mathrm{Sug}}1$,
so $Z(1)=0$. Induction on exterior degree, commuting $Z$ with
every contraction, shows that $Z(1\otimes\omega)$ is killed by
all contractions. A positive-exterior-degree element with that
property is zero, also in the indicated completion. Thus $Z=0$.
\end{proof}

Let $L$ be the degree $-1$ line derivation with
$L(B_n^a)=0$ and $L(c_{a,n})=k^{-1}\kappa^{ab}B_n^b$.
Then $[Q,L]=T$, and $X(z)=(L\otimes1)r(z)$ is the
degree-zero connection kernel of
Proposition~\ref{cylev:prop:line-kernel}.
On the twisted resolution tensor, set
\[
 \mathfrak D=\partial_z-T_{\mathcal P}\otimes1
                         +[D,L_{\mathcal P}\otimes1].
\]
Equations~\eqref{cyrec:eq:homotopy-actions} and
\eqref{cyrec:eq:homotopy-identity} give
\begin{equation}\label{cyrec:eq:connection-balance}
 \mathfrak D=\partial_z-L_1^{\mathrm{Sug}}\otimes1-R_{X(z)}.
\end{equation}
When balanced with the left-module connection
$\nabla=\partial_z+\rho_{M,N}(X)$, the last terms cancel.
The induced tensor connection is
$\partial_z-L_1^{\mathrm{Sug}}\otimes1$.

Let $p$ be a cycle of $\mathcal P$ and $\sigma$ a cycle of
$M\otimes_zN$. Suppose $\nabla\sigma$ is a boundary.
Then the derivative of $\widehat\theta(p\otimes\sigma)$ under
this affine connection is a boundary. Indeed it is
\[
 [D,L_{\mathcal P}\otimes1](\Delta_{\mathcal P}p)\otimes\sigma
                     +\Delta_{\mathcal P}p\otimes\nabla\sigma.
\]
The first term is the boundary of the homotopy applied to
$\Delta_{\mathcal P}p$, and the second is a boundary by assumption.

Here $\widehat\theta$ is obtained by including the two target
resolutions in $\widehat{\mathcal P}$.
Tensor products with $M,N$ remain algebraic and Laurent completion
is performed last. With homotopically flat modules the inclusion
of resolutions preserves the derived images.

This supplies the comparison and the translation homotopy required
at the start of Problem~\ref{cylev:prob:line-field}.
Restricted transposition can then give current Ward identities
on specified energy-graded models. The full intertwining Jacobi
identity and invertibility for an independently defined affine
fusion remain additional requirements.

\section{Two line bundles on the conifold}

An endomorphism algebra reconstructs a derived category when its
object detects every complex. The exceptional locus makes that
condition visible. On the small resolution of the conifold,
the structure sheaf misses a nonzero sheaf on the exceptional
curve. One more line bundle detects it and completes the
generating object. This is the basic conifold example of
Van den Bergh's reconstruction \cite[Introduction]{VdBFlops}.

Set
\begin{equation}\label{cyrec:eq:conifold}
 C=\mathbb C[X,Y,U,V]/(XY-UV),\qquad I=(X,U)\subset C.
\end{equation}
The gradient of the equation vanishes only at the origin.
The ring is a domain: substituting $U=qX$, $Y=qV$ embeds it in
$\mathbb C[q,X,V]$. To check injectivity, reduce monomials by
$XY=UV$ until $X$ and $Y$ do not occur together.
The image of $X^bY^dU^eV^f$, with $bd=0$, has exponents
$(d+e,b+e,d+f)$ in $(q,X,V)$, which determine the original
exponents.

The ideal $I$ has rank one and is locally free away from the
origin. At the origin its fiber $I/\mathfrak m I$ has independent
classes $X,U$, so it is not locally free there. Record their ratio
by the incidence variety
\begin{equation}\label{cyrec:eq:incidence}
 \mathcal Y=\{(X,Y,U,V,[s:t]):Xt=Us,\ Vt=Ys\}
                \subset\operatorname{Spec}C\times\mathbb P^1.
\end{equation}
Let $\rho$ forget $[s:t]$ and $\pi$ retain it.
The two charts have coordinates $(q,X,V)$ and $(q',U,Y)$,
with transitions
\begin{equation}\label{cyrec:eq:conifold-charts}
 q'=q^{-1},\qquad U=qX,\qquad Y=qV.
\end{equation}
Thus $\mathcal Y$ is the total space of
$\mathcal O_{\mathbb P^1}(-1)^{\oplus2}$ and is smooth.

Away from the origin, a nonzero row of
$\left(\begin{smallmatrix}X&U\\V&Y\end{smallmatrix}\right)$
determines $[s:t]$. Over the origin the fiber is $\mathbb P^1$.
The projection is projective, and the exceptional locus has
codimension two, so this is a small resolution.
The Jacobian of \eqref{cyrec:eq:conifold-charts} has determinant
$-1$. Consequently $dq\wedge dX\wedge dV$ and
$-dq'\wedge dU\wedge dY$ define a nowhere-zero global
three-form.

The image of $\rho^*I$ in $\mathcal O_{\mathcal Y}$ is a line
bundle $\mathcal L$. Its frames are the nonzero divisors $X$
and $U$ on the two charts, with transition $\ell_1=q\ell_0$.
Hence
\[
 \mathcal L=\pi^*\mathcal O_{\mathbb P^1}(1),
 \qquad\mathcal T=\mathcal O_{\mathcal Y}\oplus\mathcal L.
\]
The vanishing of a frame as a function on the exceptional curve
does not make its principal ideal a nonfree module.

\subsection{The full composition algebra}

A monomial $q^aX^bV^c$ extends across the second chart exactly
when $0\leq a\leq b+c$. These global functions generate $C$.
For sections of $\mathcal L$ the upper bound becomes $b+c+1$,
and multiplication by its frame $X$ identifies the sections
with $I$. For $\mathcal L^{-1}$ it becomes $b+c-1$, and division
by $X$ identifies the sections with $C+(V/X)C$.

Every map $I\to C$ is multiplication by a fraction after passing
to the fraction field. If $rX,rU\in C$, write $a=rX$.
Modulo $X$, the relation $aU=0$ lies in
$\mathbb C[Y,U,V]/(UV)$, where the annihilator of $U$ is $(V)$.
It follows that $r=b+(V/X)c$.
Both summands really map $I$ to $C$, so
\[
 I^\vee=\operatorname{Hom}_C(I,C)=C+(V/X)C.
\]

If the fraction also preserves $I$, then $Vc\in I$.
The quotient $C/I=\mathbb C[Y,V]$ has no $V$-torsion, so
$c=Xd+Ue$ and $r=b+Vd+Ye\in C$.
Therefore $\operatorname{End}_C(I)=C$ and
\begin{equation}\label{cyrec:eq:conifold-end}
 A_{\mathrm{con}}=\operatorname{End}_{\mathcal Y}(\mathcal T)
     =\begin{pmatrix}C&I^\vee\\ I&C\end{pmatrix}
     =\operatorname{End}_C(C\oplus I).
\end{equation}
The multiplication is matrix composition with the indicated
fractional entries, each of which is an actual module map.

For example, put
\begin{equation}\label{cyrec:eq:conifold-matrices}
 Z=\begin{pmatrix}0&1\\X&0\end{pmatrix},\qquad
 P_0=\begin{pmatrix}0&V/X\\U&0\end{pmatrix},\qquad
 G=\begin{pmatrix}1&0\\0&-1\end{pmatrix}.
\end{equation}
Then
\[
 Z^2=X1,\quad P_0^2=Y1,\quad
 ZP_0=\operatorname{diag}(U,V),\quad
 P_0Z=\operatorname{diag}(V,U),\quad
 [P_0,Z]=(V-U)G.
\]
These maps remain regular on $\mathcal Y$, including its
exceptional curve. A central element of $A_{\mathrm{con}}$
commutes with the two projections, hence is diagonal. Commuting
with $I\to C$ makes its diagonal entries equal.
Thus $Z(A_{\mathrm{con}})=C$.

The same algebra has a potential presentation. Let $a,c:1\to0$
and $b,d:0\to1$ be four arrows, with paths composed from right
to left, and put $\Phi=abcd-adcb$.
The Jacobi algebra is the path algebra modulo the four cyclic
derivatives of this potential.

\begin{proposition}\label{cyrec:prop:conifold-jacobi}
The Jacobi algebra of $(Q,\Phi)$ is $A_{\mathrm{con}}$.
The arrows are the four off-diagonal parts of $Z,P_0$ in
\eqref{cyrec:eq:conifold-matrices}, and its central element
\begin{equation}\label{cyrec:eq:conifold-parameter}
 (cb-ad)+(bc-da)
\end{equation}
maps to $\lambda=V-U$.
\end{proposition}

\begin{proof}
Consider
\[
 \mathcal B=\mathbb C[\lambda]\langle z,p,g\rangle/
       (g^2-1,\ gz+zg,\ gp+pg,\ pz-zp-\lambda g),
\]
with central $\lambda$. Its ordered basis is
$z^ip^jg^\epsilon$, $i,j\geq0$, $\epsilon\in\{0,1\}$.
The mixed ambiguity $gpz$ reduces in both ways to
$zpg+\lambda$, and the ambiguities with $g^2$ agree as well.
With $e_0=(1+g)/2$, $e_1=(1-g)/2$, the arrows are
$e_0ze_1,e_1ze_0,e_0pe_1,e_1pe_0$.

The cyclic derivatives are
$bcd-dcb$, $cda-adc$, $dab-bad$ and $abc-cba$.
They make \eqref{cyrec:eq:conifold-parameter} central.
The maps $z\mapsto a+b$, $p\mapsto c+d$ and
$g\mapsto e_0-e_1$ then give inverse maps between the
Jacobi algebra and $\mathcal B$.

The matrices define a surjection $\mathcal B\to A_{\mathrm{con}}$:
their arrows generate $I,I^\vee$ and the central diagonal ring.
In $\mathcal B$, the elements
\[
 X=z^2,\quad Y=p^2,\quad W=zp+\lambda g/2,
 \quad U=W-\lambda/2,\quad V=W+\lambda/2
\]
are central and obey $XY=UV$.

Filter $\mathcal B$ by the number of $z,p$ factors. Then
$\operatorname{gr}\mathcal B=\mathbb C[\lambda,z,p]\rtimes\mu_2$.
Its central subring
\[
 \mathbb C[\lambda,X,Y,W]/(XY-W^2)
\]
embeds, so its nonzero elements act injectively. The filtered algebra is therefore
torsion-free over $C$.
Over $\operatorname{Frac}(C)$, the elements $1,z,g,zg$ are
a basis with independent matrix images. The kernel of the
surjection is consequently both torsion and torsion-free, hence zero.
\end{proof}

\subsection{Generation and the inverse functor}

A compact object is one whose maps commute with arbitrary direct
sums. A generator detects zero objects by its derived Hom complex.
A compact generating bundle is \emph{tilting} when its positive
self-extensions vanish. These are properties of the entire derived
category, including objects supported on the exceptional curve.

\begin{theorem}\label{cyrec:thm:conifold-morita}
The bundle $\mathcal T$ is a compact tilting generator of
$D_{\mathrm{qc}}(\mathcal Y)$. The functors
\begin{equation}\label{cyrec:eq:conifold-morita}
 F\longmapsto R\operatorname{Hom}_{\mathcal Y}(\mathcal T,F),
 \qquad
 N\longmapsto N\otimes^{\mathbb L}_{A_{\mathrm{con}}}\mathcal T
\end{equation}
are inverse $C$-linear equivalences with
$D(\mathrm{Mod}\text{-}A_{\mathrm{con}})$.
They restrict to the categories of perfect objects.
\end{theorem}

\begin{proof}
The diagonal equation on $\mathbb P^1\times\mathbb P^1$ gives
the resolution
\[
 0\to\mathcal O(-1)\boxtimes\mathcal O(-1)
   \to\mathcal O\boxtimes\mathcal O\to\mathcal O_\Delta\to0.
\]
For any complex $F$ on $\mathbb P^1$, its integral transform
gives the triangle
\[
 \mathcal O(-1)\otimes R\Gamma(F(-1))
       \longrightarrow\mathcal O\otimes R\Gamma(F)
       \longrightarrow F.
\]
Hence $\mathcal O\oplus\mathcal O(1)$ generates.
The affine map $\pi$ has exact pushforward on quasi-coherent
sheaves. It is called \emph{conservative} because it detects zero objects.
Adjunction then proves generation by
$\mathcal T$ on $\mathcal Y$.

Moreover,
$\pi_*\mathcal O_{\mathcal Y}=
\bigoplus_{n\geq0}\mathcal O(n)^{\oplus(n+1)}$.
The self-extension groups reduce to
$H^i(\mathbb P^1,\mathcal O(n+j))$ for
$n\geq0$ and $j=-1,0,1$.
In the two-chart calculation,
\[
 H^1(\mathbb P^1,\mathcal O(d))
 =\mathbb C[q,q^{-1}]/
           (\mathbb C[q]+q^d\mathbb C[q^{-1}]),
\]
which vanishes for $d\geq-1$.
Higher cohomology vanishes because the affine cover has two terms.
The same cover and finite local freeness make derived Hom from
$\mathcal T$ commute with direct sums. This proves tilting and
compactness, also for unbounded complexes.

The right algebra action on a map $f:\mathcal T\to F$ is
$f\cdot a=f\circ a$, fixing the variance of
\eqref{cyrec:eq:conifold-morita}. Resolve a right module by
successively adjoining free generators. The tensor--Hom adjunction
unit is an equivalence on the free module $A_{\mathrm{con}}$,
by the self-extension calculation. Both functors preserve cones,
direct sums and the homotopy colimits of these free stages, so
the unit is an equivalence on every module.
The image of the counit under derived Hom is then an equivalence.
Its cone vanishes by generation, proving the other inverse identity.

Perfect complexes are the compact objects on this quasi-compact
separated scheme \cite[Tag~09M1]{StacksCompactPerfect}, and
perfect modules are the compact objects on the algebra side
\cite[Corollary~3.7]{KellerDGCategories}. Thus the equivalence
restricts to perfect objects. Every construction retains the
central $C$-action. The argument is the compact-generator form
of derived Morita reconstruction
\cite[Theorem~3.8]{KellerDGCategories}.
\end{proof}

For the exceptional inclusion $i:\mathbb P^1\hookrightarrow\mathcal Y$,
the nonzero sheaf $F=i_*\mathcal O(-1)$ has
\begin{equation}\label{cyrec:eq:exceptional-detector}
 R\operatorname{Hom}(\mathcal O,F)=0,
 \qquad
 H^1R\operatorname{Hom}(\mathcal L,F)
                 =H^1(\mathbb P^1,\mathcal O(-2))=\mathbb C.
\end{equation}
The latter class is represented by $q^{-1}$ in the Laurent
quotient. Thus the second summand is essential to the equivalence.
The ordinary center $C$ alone does not record this generating
test.

\begin{corollary}\label{cyrec:cor:conifold-fibres}
For every $c\in\mathbb C$, the restriction $\mathcal T_c$ to
$\mathcal Y_c=\{V-U=c\}$ is a compact tilting generator, with
endomorphism algebra $A_{\mathrm{con}}/(\lambda-c)A_{\mathrm{con}}$.
The equivalence \eqref{cyrec:eq:conifold-morita} holds on each
fiber with this quotient algebra.
\end{corollary}

\begin{proof}
Multiplication by $\lambda-c$ is injective on
$\mathcal E nd(\mathcal T)$, since it is a coordinate on each
affine chart. Its restriction sequence and the preceding
cohomology vanishing compute the stated endomorphism algebra
and zero positive self-extensions. Closed pushforward from the
fiber is conservative, and its adjunction with $\mathcal T$
proves generation. The same affine cover proves compactness.
The tensor--Hom proof therefore applies on every fiber.
\end{proof}

The fiber charts have transition
$q'=q^{-1}$, $Y=q^2X+cq$. They are smooth for every $c$.
At $c=0$ the fiber is $T^*\mathbb P^1$.
For $c\ne0$, it is the smooth affine surface
$XY=U(U+c)$, which avoids the conifold origin.
This parameter describes surfaces inside the fixed threefold
$XY=UV$. It is distinct from a smoothing parameter in an
equation $XY-UV=c$.

\section{Intrinsic Hall fibers and ordinary cooperation}

The derived equivalence of the conifold retains modules and their
extensions. A cooperation on the Hall coefficients asks for more:
it must preserve the extension product in a specified tensor
category. Specializing a coefficient map to polynomials can obscure
that question. Its kernel can grow even when the coefficient module
is free over the parameter ring. The conifold gives an obstruction
supported on precisely these surviving classes.

Let the torus $(\mathbb C^*)^2$ act on the four arrows
$(a,b,c,d)$ with weights $(\alpha,\beta,-\alpha,-\beta)$.
Put
\[
 R=\mathbb C[\alpha,\beta],\qquad
 P=\alpha^2,\quad Q=\beta^2,\quad\delta=P-Q,
\]
with $|\alpha|=|\beta|=2$. These equivariant parameters are
independent of the central function $\lambda=V-U$ above.

For charge $\mathbf d=(n,m)$ use the representation stack
$\mathfrak M_{n,m}$ of the four-arrow quiver, the function
$f_{n,m}=\operatorname{Tr}\Phi$, and the critical complex
\begin{equation}\label{cyrec:eq:critical-carrier}
 H_{n,m}=R\Gamma_T\bigl(\mathfrak M_{n,m},
                  \phi^p_{f_{n,m}}\mathbb C[4nm-n^2-m^2]\bigr).
\end{equation}
Here $\phi^p_f=\phi_f[-1]$, so $\phi^p_0\mathbb C=\mathbb C$.
Use the critical-chart orientation with ordered arrow and vertex
determinants, and twist the product of charges $(n,m),(r,s)$
by $(-1)^{mr}$. All sums in charge and degree are direct sums.

The critical transformations on invariant flags give an augmented
unital associative algebra on these complexes. Their hypotheses
hold because all representations are allowed, the potential is
homogeneous with every critical value zero, and the Euler form
\[
 \chi((n,m),(r,s))=nr+ms-2ns-2mr
\]
is symmetric. The product therefore has degree zero
\cite[\S3.2]{DavisonCriticalCoHA}.
Chains of flag refinements give its coherent higher
identifications.

Choose a charge-graded K-flat dg algebra $B$
representing it, with $B_{0,0}=R$.
Such a model can be obtained by bar and cobar from a strictly
unital semifree higher-algebra representative. Reduced word length
in charge $(n,m)$ is at most $n+m$, so this requires no charge
completion and assumes no formality.

\subsection{Integral coefficient modules}

Deleting the arrow $d$ writes the potential as
$\operatorname{Tr}(d(abc-cba))$. Dimensional reduction identifies
the critical coefficients with shifted equivariant Borel--Moore
homology of
\[
 Z_{n,m}=\{abc-cba=0\}
\]
in the three-arrow representation space. The shift is
$2nm+\chi((n,m),(n,m))$
\cite[Theorem~A.1 and Corollary~A.9]{DavisonCriticalCoHA}.

Closed pushforward gives a coefficient map
\begin{equation}\label{cyrec:eq:hall-coefficient}
 \theta_{n,m}:H^*(B_{n,m})\longrightarrow
 R[x_1,\ldots,x_n,y_1,\ldots,y_m]^{\mathfrak S_n\times\mathfrak S_m}
                   [-\chi((n,m),(n,m))].
\end{equation}
The roots and torus parameters have raw weight one and
cohomological degree two. An output root polynomial of raw
weight $w$ has Hall degree $2w+\chi((n,m),(n,m))$.
Position coordinates will be introduced separately.

In charges $(1,0)$, $(0,1)$ and $(1,1)$ the maps identify
\begin{equation}\label{cyrec:eq:small-coefficients}
 H^*(B_{1,0})=R[x][-1],\quad
 H^*(B_{0,1})=R[y][-1],\quad
 H^*(B_{1,1})=R[x,u][2],\qquad u=y-x.
\end{equation}
In the mixed charge all matrices are scalar and the potential
vanishes. If $a_r=x^r$ and $b_s=y^s$ are the simple-charge
states, the two flags give
\begin{equation}\label{cyrec:eq:mixed-hall-product}
 a_rb_s=x^ry^s(u^2-Q),\qquad
 b_sa_r=-x^ry^s(u^2-P).
\end{equation}
The first flag sets $b=d=0$ and the second sets $a=c=0$.
Their normal Euler classes are the two quadratics, with the
second sign supplied by the chosen twist.
Thus their products span the proper ideal
$(u^2-P,\delta)$, while the full mixed coefficient contains
the degree $-2$ class $1_{(1,1)}$.

For charge $(2,1)$ put
\[
 L=x_1+x_2-2y,\qquad H=(x_1-y)(x_2-y)+P.
\]
For $(1,2)$ use
$L'=y_1+y_2-2x$, $H'=(y_1-x)(y_2-x)+P$.
Then the coefficient maps are injective with images
\begin{equation}\label{cyrec:eq:next-ideals}
 I_{21}=(L,H,\delta)R[y,L,H],\qquad
 I_{12}=(L',H',\delta)R[x,L',H'],
\end{equation}
both with Hall shift $[3]$.
In charge $(2,1)$ the arrows $a,c$ are columns, and the cut equation is
\[
 \det[a\ c]\,(b_2,-b_1)^{\mathsf t}=0.
\]
Its complement has $[a\ c]$ invertible and $b\ne0$.

Gauge the two columns to a basis. The remaining punctured vector
$(ba,bc)$ has weights $\beta+\alpha,\beta-\alpha$ and Euler
class $-\delta$. Its equivariant cohomology is $R[y]/(\delta)$
and is even. Restriction from the ambient ring is surjective,
setting $L=H=0$. The closed--open Borel--Moore sequence is
therefore injective on the closed locus, with kernel ideal
$I_{21}$. The row version of the same argument proves the
formula for $I_{12}$.

The next output charge requires one more coefficient. Put
\begin{equation}\label{cyrec:eq:two-two-coordinates}
\begin{gathered}
 \tau=y_1+y_2,\qquad \xi=x_1+x_2-y_1-y_2,\\
 A_0=(x_1-x_2)^2,\quad C_0=(y_1-y_2)^2,\quad\zeta=A_0-C_0,\\
 h=\frac{\xi^2-A_0-C_0}{2}+P+Q,\qquad g=\delta\xi,\\
 k_0=\frac{\zeta^2-\xi^4+(12P-4Q)\xi^2-16PQ}{16}.
\end{gathered}
\end{equation}
The invariant coefficient ring is
$\mathscr B=R[\tau,\xi,h,\zeta]$.

\begin{lemma}\label{cyrec:lem:two-two-ideal}
The map $\theta_{2,2}$ is injective with image
$I_{22}=(h,k_0,g)\subset\mathscr B$ and Hall shift $[8]$.
The sequence $h,k_0,g$ is regular.
\end{lemma}

\begin{proof}
For $2\times2$ matrices, expansion in the entries of $d$ gives
\[
 \operatorname{Tr}(abcd-adcb)
 =\det[\operatorname{vec}(a),\operatorname{vec}(c),
               \operatorname{vec}(\operatorname{adj}b),
               \operatorname{vec}(\operatorname{adj}d)].
\]
The cut locus is the rank-at-most-two locus of the first three
columns, a generic $4\times3$ matrix. Its complement is the
bundle of injections from a rank-three bundle $U_3$ into a
rank-four bundle $E_4$. Here the roots of $E_4$ are $x_i-y_j$,
and those of $U_3$ are $-\alpha,\alpha,\xi-\beta$.
Division of their Chern polynomials gives
\begin{align*}
 c_2(E_4-U_3)&=h,\\
 c_3(E_4-U_3)&=\beta h+g,\\
 c_4(E_4-U_3)&=k_0+\tfrac14\xi^2h
                   +(P-\beta\xi+Q)h-(\xi-\beta)g.
\end{align*}

The universal injection bundle has reductive stabilizer
$GL_3\times GL_1$. Its characteristic-class restriction sends
$E_4$ to $U_3\oplus L$ and has kernel generated by
$c_2,c_3,c_4$ of the difference bundle.
Pulling back its polynomial cochain model gives the Koszul
complex on these three classes. The classifying bases are simply
connected, so this computes the equivariant cohomology degreewise.

In the present ring, $h$ is a coordinate and $16k_0$ is monic
in $\zeta$ of degree two. Modulo these two elements the ring is
free on $1,\zeta$ over $R[\tau,\xi]$. Multiplication by
$g=\delta\xi$ is injective on that free module, proving
regularity. Hence the complement has even cohomology
$\mathscr B/I_{22}$, and restriction is surjective.
The closed--open sequence makes the coefficient map injective
with image $I_{22}$. Equation~\eqref{cyrec:eq:hall-coefficient}
gives the shift.
\end{proof}

The critical flag product is integral and intertwines these
coefficient maps with ambient flag induction. This follows from
the general-cut comparison, which retains the full off-diagonal
derivative $abc-cba$ in its Cartesian flag square
\cite[Appendix~A, Theorem~A.2]{YangZhaoGeneralCut}.
In the ambient composite, the Euler class of the diagonal
inclusion cancels its inverse in the refined pullback formula.
The remaining Euler classes are exactly the four forbidden arrow
blocks. Proper parabolic induction contributes the same-color
Grassmannian denominators.

Explicitly, for a splitting $(n,m)+(r,s)$, let $L_0,L_1$ be
subsets of the output roots of sizes $n,m$, and let $M_0,M_1$
be their complements. The first factor is the subobject.
The product of coefficient polynomials is
\begin{equation}\label{cyrec:eq:flag-formula}
 (-1)^{mr}\sum_{L_0,L_1}
 \frac{f(L_0,L_1)g(M_0,M_1)
  \displaystyle\prod_{u\in L_0,v\in M_1}((v-u)^2-Q)
  \displaystyle\prod_{u\in L_1,v\in M_0}((v-u)^2-P)}
 {\displaystyle\prod_{u\in L_0,v\in M_0}(v-u)
  \displaystyle\prod_{u\in L_1,v\in M_1}(v-u)}.
\end{equation}

Each term is one fixed flag, so there is no factorial.
Proper induction makes the sum polynomial. Equality with the
critical product follows first after localization, then integrally
because the ambient polynomial ring embeds in that localization.
In each injective output charge above, this determines the actual
critical class, including products of full coefficients outside
the simple-current subalgebra.

\subsection{Classes that survive specialization}

Let $S=R/(\delta)$ and take the intrinsic derived coefficient
fiber $B_S=B\otimes_R^{\mathbb L}S$. The base ring has the
resolution
\[
 K_\delta=R\otimes\Lambda(\epsilon),\qquad
 |\epsilon|=3,\quad d\epsilon=\delta.
\]
Its odd generator has internal degree four and resolution degree
$-1$. The K-flat model $B\otimes_RS$ represents the same fiber
as an $S$-algebra. This is coefficient base change. Identifying it
with critical cohomology for a changed geometric equivariance
would require a separate theorem.

The ideals in \eqref{cyrec:eq:next-ideals} are free over $R$,
as the decompositions
\begin{equation}\label{cyrec:eq:free-mixed}
 I_{21}=(L,H)R[y,L,H]\oplus\delta R[y],\qquad
 I_{12}=(L',H')R[x,L',H']\oplus\delta R[x]
\end{equation}
show. Denote the unique classes with coefficient $\delta$ by
$\Gamma_D,\Gamma_E$. Both have Hall degree one.
Their images $D,E$ in $H^*(B_S)$ are nonzero, although polynomial
evaluation kills them. Projection to the constant term of their
second summands gives $S$-linear coefficient detectors
$D\mapsto1$, $E\mapsto1$.

Likewise monic division first by $h$ and then by $16k_0$ gives
the free $R$-module decomposition
\begin{equation}\label{cyrec:eq:free-two-two}
 I_{22}=h\mathscr B\oplus k_0R[\tau,\xi,\zeta]
                  \oplus\delta\xi R[\tau,\xi]\{1,\zeta\}.
\end{equation}
Let $\Gamma_g$ be the unique Hall class with coefficient
$g=\delta\xi$. Its Hall degree is $-2$, and its intrinsic
image $G$ also survives. These free decompositions imply that
the relevant derived tensors over $S$ are the ordinary tensors
of specialized cohomology modules.
Indeed graded-free cohomology admits a choice of lifted basis
cycles giving a quasi-isomorphism from its zero-differential
complex. Base change preserves that comparison.

In the simple and first mixed charges define marked states
\[
 a=1_{(1,0)},\quad b=1_{(0,1)},\quad
 q=1_{(1,1)},\quad q_x=x_{(1,1)},\quad
 r_0=u_{(1,1)},\quad r_1=(xu)_{(1,1)}.
\]
Their degrees are $1,1,-2,0,0,2$, respectively.
The subscript specifies the charge. In particular $q$ is not
the Hall unit, which belongs to charge zero.

\begin{proposition}\label{cyrec:prop:intrinsic-relations}
The actual critical Hall classes over $R$ satisfy
\begin{equation}\label{cyrec:eq:intrinsic-relations}
 [a,r_0]=-\Gamma_D,\quad [b,r_0]=\Gamma_E,\quad
 [q,q_x]=-2\Gamma_g,\quad
 [r_0,r_1]+\Gamma_Ea+2P\Gamma_g=0.
\end{equation}
They retain these identities under intrinsic base change to
$S$ and to the homogeneous origin $\alpha=\beta=0$.
\end{proposition}

\begin{proof}
All displayed commutators have an even input.
The two flags in charge $(2,1)$ give
\[
 \theta_{2,1}([a,f])
       =\delta\frac{f(x_2,y)-f(x_1,y)}{x_2-x_1}.
\]
For $f=y-x$ this is $-\delta$. The reflected divided difference
is $+1$, giving $+\delta$ in charge $(1,2)$.

For two mixed inputs denote the output roots by $X,U;Y,V$,
and put
\[
 A_1=Y-X,\quad B_1=V-U,\quad C_1=V-X,\quad D_1=Y-U,
 \qquad\xi=X+U-Y-V.
\]
Pairing complementary flags uses
\[
 (D_1^2-Q)(C_1^2-P)-(C_1^2-Q)(D_1^2-P)
                  =\delta(C_1^2-D_1^2).
\]
The two pairs in the commutator $[q,q_x]$ then give
$-2\delta\xi$, proving the third relation by injectivity.

For $[r_0,r_1]$ the remaining coefficient is
\[
 \frac{\delta}{V-Y}
       \bigl(A_1B_1(C_1^2-D_1^2)+C_1D_1(B_1^2-A_1^2)\bigr)
       =-\delta\xi K,
\]
where $K=X^2+U^2-(X+U)(Y+V)+2YV$.
The two flags for $\Gamma_Ea$ instead give
$\delta(F(U)-F(X))/(U-X)=\delta\xi(K-2P)$, where
$F(t)=((t-Y)^2-P)((t-V)^2-P)$.
Adding $2P\delta\xi$ gives zero.
Injectivity before specialization proves the relation on Hall
classes. The free models above then transport it to the intrinsic
fibers, including the classes whose polynomial images vanish.
\end{proof}

\begin{theorem}\label{cyrec:thm:hall-obstruction}
There is no degree-zero, charge-preserving unital map
\[
 B_S\longrightarrow B_S\otimes_S^{\mathbb L}B_S
\]
preserving binary multiplication in $D(S)$, with the ordinary
Koszul interchange and charge-zero projection as a two-sided
counit. The same conclusion holds at the intrinsic homogeneous
origin. Coassociativity is not needed for this obstruction.
\end{theorem}

\begin{proof}
Write $\overline\Delta(v)=\Delta(v)-v\otimes1-1\otimes v$.
The counits fix the charge-zero summands. The simple charges
have no proper positive splits, so $a,b$ are primitive.
The only proper splits of charge $(1,1)$ have minimum degree
two, whereas $q,q_x,r_0$ have degrees $-2,0,0$.
They too are primitive.

In the ordinary graded tensor product,
a supercommutator of primitives is primitive.
The first three relations in \eqref{cyrec:eq:intrinsic-relations}
therefore force $D,E,G$ to be primitive.

The full possible correction to the degree-two state $r_1$ is
\[
 \overline\Delta(r_1)=\lambda a\otimes b+\mu b\otimes a,
 \qquad\lambda,\mu\in\mathbb C.
\]
Since $r_0$ is even and $a,E$ are odd, the last relation then
has reduced coproduct
\begin{equation}\label{cyrec:eq:four-defects}
 \lambda D\otimes b+\mu b\otimes D
                -(\lambda+1)a\otimes E+(1-\mu)E\otimes a.
\end{equation}

The four summands have different charge splits. Their coefficient
detectors take value one, by \eqref{cyrec:eq:free-mixed} and
the constant simple-charge coefficients. The first forces
$\lambda=0$, while the third forces $\lambda=-1$.
The same free decompositions and detectors survive at the origin.
The contradiction occurs in cohomology of the actual derived
tensor components, so no choice of chain homotopy removes it.
\end{proof}

The four equations can be expressed by the two-term complex
\begin{equation}\label{cyrec:eq:relative-obstruction}
 S^2\xrightarrow{\Psi}S^4,
 \qquad
 \Psi=\begin{pmatrix}1&0\\0&1\\-1&0\\0&-1\end{pmatrix},
 \qquad\omega=(0,0,1,-1)^{\mathsf t}.
\end{equation}
The map $J(v_1,v_2,v_3,v_4)=(v_1+v_3,v_2+v_4)$ has
kernel $\operatorname{im}\Psi$ and a section
$(v,w)\mapsto(0,0,v,w)$. Hence the obstruction class maps to
$J\omega=(1,-1)$ in its cokernel $S^2$.
Even allowing arbitrary $S$-valued corrections cannot kill it.
This is a quotient of the necessary multiplication equations,
without an identification with the full bialgebra deformation
complex.

Deleting the two terms containing $D$ would give the apparent
solution $(\lambda,\mu)=(-1,1)$. Its residual on the intrinsic
fiber is $-D\otimes b+b\otimes D\ne0$.
Thus passing to the specialized polynomial image changes the
coefficient object on which the obstruction is tested.

\subsection{The regular collision tensor}

The coefficient tensor in the theorem is forced by a specific
regular collision construction. Let the position curve be
$X=\mathbb A^1$ with coordinate $z$, independent of all Chern
roots.

For a partition $\pi$ of $[N]$, a run is a maximal
consecutive interval with the same block label, and $\ell(\pi)$
is the number of runs. In particular $13|2$ has three runs.
Let $i_\pi$ repeat block coordinates and $j_\pi$ include their
pairwise-distinct locus. With trivial coefficient connection put
\[
 M_\pi(B)=i_{\pi,*}j_{\pi,*}j_\pi^!
            (\omega_{X^{|\pi|}}\otimes B^{\otimes_R\ell(\pi)}).
\]
The canonical module here is unshifted.

A merger multiplies exactly the adjacent runs that acquire the
same label. It leaves all three factors on $13|2$ and multiplies
them in order on the subsequent total collision.
Order blocks by their largest elements and use the full normal
transfer $v^{-r-1}dv\leftrightarrow\delta_v\partial_v^r/r!$.
The run complex is
\begin{equation}\label{cyrec:eq:run-complex}
 C_0=R,\qquad C_N=\bigoplus_{\pi\vdash[N]}M_\pi(B)[|\pi|],
 \qquad
 d|_\pi=(-1)^{|\pi|}d_B+\sum_{i<j}(-1)^{i-1}r_{ij}.
\end{equation}
Here $d_B$ means the tensor differential on all run factors.
The principal-part map $r_{ij}$ includes both the normal transfer
and the ordered multiplication of the newly adjacent factors.

For disjoint mergers the state maps agree and the signs cancel.
For three overlapping blocks the state map is the same ordered
product on each route. The scalar identity is $L-R+N=0$ for
the three full principal-part transfers. In coordinates
$s=z_1-z_3$, $t=z_2-z_3$, its change of normal coordinates is
\[
 (s-t)^{-a}t^{-b}\longmapsto
     \sum_{j=0}^{b-1}\binom{a+j-1}{j}s^{-a-j}t^{-b+j}.
\]
Partial fractions at $s=0$ and $s=t$ prove the identity, and
right differential linearity retains all normal derivatives.
Thus \eqref{cyrec:eq:run-complex} has square-zero differential.

On a cut open, every block crossing that cut restricts to zero.
The remaining partitions split uniquely, and their runs concatenate.
This gives a strong cut coalgebra. Extraordinary pullback along a
monotone surjection retains exactly the partitions pulled back
from its target and deletes repeated adjacent run labels.
These identifications compose and preserve the residues.

The cut bar--cobar equivalence of Theorem~\ref{cycuts:thm:duality}
then gives its cut cobar $A_\Phi$.
The gap contraction retains exactly partitions crossing every
cut, with coefficient $M_\pi(B)[|\pi|-1]$.
The finite perturbation
$\sum_{a=0}^{N-1}(-\partial_{\mathrm{col}}h_0)^a$,
where $h_0$ contracts the first admissible gap, supplies the
comparison and its higher operations.

The scalar unit of $B$ supplies the local vacuum. Its $N$th coefficient is
\[
 \boldsymbol\omega_{X^N}\otimes R,
 \qquad \boldsymbol\omega_Y=\omega_Y[\dim_{\mathbb C}Y].
\]
It is distinct from the empty cut unit, which has only arity zero.
Under coefficient base change, the finite partition sums, residue
maps and gap contractions commute with derived scalar extension.
Thus the same construction represents $B_S$.

The decisive tensor calculation is already binary. On a normal
complex disk $D$, let $j:D^*\hookrightarrow D$ and
$i:\{0\}\hookrightarrow D$. For a map of coefficient complexes
$\mu:V\to W$, define the regular attachment
\begin{equation}\label{cyrec:eq:regular-attachment}
 G(\mu)=\operatorname{Fib}\left(
 Rj_*\boldsymbol\omega_{D^*}\otimes V
      \xrightarrow{\operatorname{Res}\otimes\mu}i_*W[1]\right).
\end{equation}
Its open coefficient is $V$ and its costalk is $W$.
After removing the common tangent density, the binary run
coefficient is $G(\mu_B:B\otimes_R^{\mathbb L}B\to B)$.

\begin{lemma}\label{cyrec:lem:attachment-tensor}
For perfect coefficient arrows, and for their Ind-extension with
K-flat coefficients,
\[
 G(\mu)\otimes^!G(\nu)\simeq G(\mu\otimes^{\mathbb L}\nu).
\]
A map of attachments inducing open and closed maps $a,b$ obeys
$b\mu=\nu a$ in the derived category.
\end{lemma}

\begin{proof}
For perfect coefficients, Verdier duality identifies $G(\mu)$
with the dual of a constructible diagram. Its closed stalk is
$W^\vee$, its punctured stalk is $V^\vee$, its monodromy is the
identity and its generization is $\mu^\vee$. Ordinary tensor tensors its stalks
and generization, and duality gives the exceptional tensor formula.
Every coefficient arrow is a filtered homotopy colimit of perfect
arrows. The fixed two-stratum kernels, their localization triangles
and derived tensors commute with these colimits, so the formula
extends without dualizing nondualizable coefficients.
Naturality of the localization boundary then gives $b\mu=\nu a$
on its normal-circle residue component.
\end{proof}

Transport the exceptional tensor through cut bar and cobar:
\[
 A\diamond A'=\Omega_{\mathrm{cut}}
       (B_{\mathrm{cut}}A\otimes^!B_{\mathrm{cut}}A').
\]
Suppose a unital cooperation $A_\Phi\to A_\Phi\diamond A_\Phi$
preserves cuts, collisions and the local vacuum, with a two-sided
counit to the transported vacuum unit. Its unary coefficient is
a map $\Delta_B:B\to B\otimes_R^{\mathbb L}B$.
The cut comparison forces its open binary map to be two copies
of $\Delta_B$, with the middle factors interchanged. Its closed
map is $\Delta_B$ by the common-root comparison.

Lemma~\ref{cyrec:lem:attachment-tensor} therefore forces
\[
 \Delta_B\mu_B=(\mu_B\otimes\mu_B)
            (1\otimes\tau\otimes1)(\Delta_B\otimes\Delta_B)
\]
with the ordinary Koszul swap $\tau$. The geometric counit
induces the charge-zero projection.
Theorem~\ref{cyrec:thm:hall-obstruction} excludes this cooperation
on the intrinsic fiber $S$ and at the homogeneous origin.

Over $R$ itself there is an earlier obstruction. The mixed
products \eqref{cyrec:eq:mixed-hall-product} satisfy
\[
 \{a_{r+2},b_s\}-2\{a_{r+1},b_{s+1}\}
       +\{a_r,b_{s+2}\}-P a_rb_s-Q b_sa_r=0.
\]
Both simple charges are odd and primitive in any such cooperation.
Projection of the proposed coproduct onto charge split
$(1,0)|(0,1)$ gives $-\delta a_r\otimes b_s\ne0$.
This class remains nonzero after inverting $\delta$.
Its vanishing on $S$ is the reason the surviving intrinsic
classes were necessary for the stronger fiber calculation.

A tensor with meromorphic interchange or different collision
attachment changes these hypotheses. The rational Hall tensor
of the preceding chapter is an example of a construction with a
specified changed interchange. Comparing it with the global
critical coefficients requires a map retaining the intrinsic
classes, the vacuum, the counit and the triple-collision identities.

\section{Derived interfaces at a node}

There is an associative tensor available without a coproduct on
the coefficient algebra: composition of bimodules by derived
tensor over that algebra. At a singular point its higher Tor
classes are essential. The node gives both an explicit composition
and a test of how much of the center one branch can detect.

Let $k$ have characteristic zero and put
\[
 A=k[u,m]/(um),\qquad P=A/(m),\qquad Q=A/(u),
 \qquad A^e=A\otimes_kA^{\mathrm{op}}.
\]
Regard the two branch modules as diagonal $A$-bimodules.
In $D(A^e)$, interface composition is
$M\otimes_A^{\mathbb L}N$, with unit $A$.
The two-sided bar complex gives its associator: both
parenthesizations realize the same multibar complex, and its
successive face maps give the coherence identities.
All tensors here are algebraic.

\subsection{The derived intersection retains infinitely many classes}

A semifree resolution of $P$ as an $A$-algebra is
\begin{equation}\label{cyrec:eq:node-resolution}
 R_P=A[e,t],\qquad |e|=-1,\quad |t|=-2,
 \qquad de=m,\quad dt=ue.
\end{equation}
The odd variable is exterior. Use divided powers
$t^{[j]}=t^j/j!$ for its even variable.
The underlying free complex has one cell in each nonpositive
degree and differential alternately $m,u,m,u,\ldots$.
Since $\operatorname{ann}(m)=uA$ and
$\operatorname{ann}(u)=mA$, it is exact away from degree zero,
where its cohomology is $P$.
The analogous resolution of $Q$ interchanges $u$ and $m$.

\begin{proposition}\label{cyrec:prop:node-composition}
There are quasi-isomorphisms of derived $A$-bimodules
\begin{align}
 P\otimes_A^{\mathbb L}Q&\simeq k[t],
               &|t|&=-2,\label{cyrec:eq:node-binary}\\
 P\otimes_A^{\mathbb L}Q\otimes_A^{\mathbb L}P
      &\simeq k[t_1,t_2]\otimes\Lambda(\epsilon),
               &|\epsilon|&=-1.
 \label{cyrec:eq:node-triple}
\end{align}
The targets have zero differential, and both $A$-actions factor
through $A/(u,m)=k$.
\end{proposition}

\begin{proof}
Tensor $R_P$ with $Q$. This gives $k[m,e,t]$ with $de=m$,
$dt=0$. Projection killing $m,e$ is a bimodule map to $k[t]$.
The regular Koszul pair contracts by
$H(m^j)=em^{j-1}$ for $j>0$, with $H$ zero on the constant
and odd terms. The homotopy is $k$-linear, which suffices to
prove that the bimodule projection is a quasi-isomorphism.

For the triple tensor, resolve both outer copies of $P$.
The common complex is
\[
 k[m,e_1,e_2,t_1,t_2],\qquad
 de_1=de_2=m,\quad dt_1=dt_2=0.
\]
Put $\epsilon=e_2-e_1$ and contract $(m,e_1)$.
The remaining variables give \eqref{cyrec:eq:node-triple}.
Flatness of the outer resolutions makes this one complex compute
both parenthesizations.
\end{proof}

The ordinary intersection $P\otimes_AQ=k$ discards all positive
powers of $t$. Replacing the first derived intersection by that
point changes its next composition to
$k\otimes_A^{\mathbb L}P\simeq k[t]\otimes\Lambda(e)$.
Its degree $-2$ has dimension one, whereas that degree in
\eqref{cyrec:eq:node-triple} has dimension two.
The distinction is present before any completion.

The bounded coherent category therefore does not contain the
composite, even though it contains both branches.
Nor is $\operatorname{Perf}(A^e)$ closed under this composition:
\[
 A^e\otimes_A A^e\simeq A\otimes_kA\otimes_kA
\]
is an infinite-rank free module over the two outer factors $A^e$.
Its degree-zero cohomology is not finitely generated over the
Noetherian ring $A^e$, so it cannot be perfect.

There is a different closed class. Bimodules perfect as left
$A$-modules and as right $A$-modules are closed under relative
derived tensor and contain the unit. For left perfection, express
the second factor by finite cones and retracts of left free
modules. Tensoring turns these into finite cones and retracts of
the first factor. The right-module argument is the reverse.

The branch modules do not belong to this class, as their periodic
self-extensions below show. Their composition is instead defined
in the full $D(A^e)$.

\subsection{Cross-branch extensions and their composition}

Let $P_\bullet,Q_\bullet$ be the alternating free resolutions
just constructed. Applying Hom to them gives
\begin{equation}\label{cyrec:eq:branch-ext}
\begin{aligned}
 \operatorname{Ext}^0_A(P,P)&=P,&
 \operatorname{Ext}^{2j}_A(P,P)&=k\quad(j\geq1),&
 \operatorname{Ext}^{2j+1}_A(P,P)&=0,\\
 \operatorname{Ext}^{2j+1}_A(P,Q)&=k\quad(j\geq0),&
 \operatorname{Ext}^{2j}_A(P,Q)&=0\quad(j\geq0).
\end{aligned}
\end{equation}
For the first row the cochain differentials alternate between
zero and multiplication by $u$ on $k[u]$.
For the second they alternate between multiplication by $m$
and zero on $k[m]$. The formulas with the branches exchanged
follow by symmetry.

The degree-one cross classes have nonzero Yoneda composite.
Define $F:P_\bullet\to Q_\bullet$ by sending the cell in
degree $-n$, for $n\geq1$, to $(-1)^{n-1}$ times the cell
in degree $-n+1$, and killing degree zero.
Define $G:Q_\bullet\to P_\bullet$ by the same coefficients.
Adjacent coefficients cancel, so $dF+Fd=dG+Gd=0$.
Their composites equal minus the two-step shift, which represents
the nonzero class in $P/uP=\operatorname{Ext}^2_A(P,P)$.

The periodic matrices make the sign visible:
\begin{equation}\label{cyrec:eq:node-matrix-factorizations}
 d_P=\begin{pmatrix}0&m\\u&0\end{pmatrix},\qquad
 d_Q=\begin{pmatrix}0&u\\m&0\end{pmatrix},\qquad
 F=\begin{pmatrix}0&1\\-1&0\end{pmatrix}.
\end{equation}
The two differentials square to $um$, while
$d_QF+Fd_P=0$ and $F^2=-1$.
Thus these are matrix factorizations of the nodal equation,
with the same nonzero cross-branch composition after
periodization.

\subsection{The Hochschild class detected by one branch}

The ordinary center of the commutative node is $A$.
Its derived center also sees the singular point through a
periodic degree-two class. To calculate it, use left and right
coordinates on $A^e$ and the diagonal Tate resolution
\begin{equation}\label{cyrec:eq:node-diagonal}
 \mathcal T_A=A^e\otimes\Lambda(e_u,e_m)\otimes\Gamma(h),
 \quad |e_u|=|e_m|=-1,\quad |h|=-2,
\end{equation}
with
\[
 de_u=u_L-u_R,\qquad de_m=m_L-m_R,\qquad
 dh=m_Le_u+u_Re_m.
\]
The last differential squares to $u_Lm_L-u_Rm_R=0$.
The even variable has divided powers and
$dh^{[r]}=(dh)h^{[r-1]}$.

To verify exactness, put $a=u_L-u_R$, $b=m_L-m_R$,
$x=u_R$, $y=m_R$.
Then $A^e=A[a,b]/(ay+bx+ab)$.
The defining polynomial is a nonzero divisor, with leading
monomial $ab$ and coefficient one. The Koszul complex on $a,b$
therefore has, after this quotient, homology $A$ in degrees
zero and $-1$. The latter generator is $m_Le_u+u_Re_m$.

Filtering \eqref{cyrec:eq:node-diagonal} by divided-power degree
pairs each positive divided power with this generator and
contracts them. The filtration is finite in each chain degree,
so the augmentation to $A$ is a free diagonal resolution.

Applying $\operatorname{Hom}_{A^e}(-,A)$ gives the additive
cochain model
\begin{equation}\label{cyrec:eq:node-hochschild-complex}
 A[\tau]\otimes\Lambda(\xi_u,\xi_m),
 \qquad |\tau|=2,\quad |\xi_u|=|\xi_m|=1,
 \qquad d\xi_u=m\tau,\quad d\xi_m=u\tau.
\end{equation}
The exterior notation specifies the basis and differential.
Products will be obtained from maps of the resolution, without
identifying this exterior product with a strict Hochschild cup
product.

\begin{proposition}\label{cyrec:prop:node-hochschild}
With $E_u=u\xi_u$ and $E_m=m\xi_m$, one has
\begin{align*}
 \operatorname{HH}^0(A)&=A,\\
 \operatorname{HH}^1(A)&=u k[u]\xi_u\oplus m k[m]\xi_m,\\
 \operatorname{HH}^{2r}(A)&=k\tau^r\quad(r\geq1),\\
 \operatorname{HH}^{2r+1}(A)&=kE_u\tau^r=kE_m\tau^r
                                      \quad(r\geq1).
\end{align*}
Multiplication by $\tau$ is periodic in degrees at least two.
The Gerstenhaber brackets are
\begin{equation}\label{cyrec:eq:euler-periodic}
 [E_u,\tau]=[E_m,\tau]=-\tau.
\end{equation}
The products $E_u\tau=E_m\tau$ are nonzero, while all products
of two degree-one classes vanish.
\end{proposition}

\begin{proof}
A degree-one cycle has coefficients $A_1,A_2$ with
$mA_1+uA_2=0$. Since $(u)\cap(m)=0$, they lie in $(u)$
and $(m)$, respectively. In positive even degree the scalar
$\tau$ coefficient is reduced modulo $(u,m)$, and no nonzero
coefficient of $\xi_u\xi_m$ is annihilated by both $u$ and $m$.
In odd degrees at least three, the boundary
$(m\xi_m-u\xi_u)\tau$ identifies the Euler terms and kills
their higher polynomial multiples. This computes the groups.

The chain map lowering the divided power of $h$ by one and
killing $h^{[0]}$ represents $\tau$.
It commutes with the divided-power differential, including its
first step. Precomposition with it gives cup multiplication by
$\tau$, proving periodicity and the asserted nonzero products.

Scale $u,m$ independently by a two-dimensional torus.
The resolution generator $h$ has weight $(1,1)$, so $\tau$
has weight $(-1,-1)$. Bracket with an Euler derivation is this
weight action, proving \eqref{cyrec:eq:euler-periodic}.
Degree-one classes have weights $(r,0)$ or $(0,s)$ with
$r,s\geq0$, whereas degree two has only weight $(-1,-1)$.
Weight preservation forces all degree-one pair products to vanish.
\end{proof}

A Hochschild class acts on a module by tensoring the corresponding
diagonal endomorphism with that module. For the branch $P$, this
is the characteristic map
\[
 \chi_P:\operatorname{HH}^*(A)\longrightarrow
                                      \operatorname{Ext}_A^*(P,P).
\]
Let $\beta$ denote the positive two-step shift of $P_\bullet$.
Then
\begin{equation}\label{cyrec:eq:characteristic-values}
 \chi_P(\tau)=\beta\ne0,\qquad\chi_P(E_u)=0.
\end{equation}

Indeed, after tensoring the diagonal resolution on the right with
$P$, put $h'=h+e_ue_m$. The differential becomes $dh'=u_Le_m$.
The regular pair $(u_L-u_R,e_u)$ contracts, leaving exactly
$A[e_m,h']$ with $de_m=m_L$, $dh'=u_Le_m$.
Lowering the divided power of $h$ becomes lowering that of $h'$,
which is the branch periodic shift. The degree-one image vanishes
by \eqref{cyrec:eq:branch-ext}.

Consequently no $E_2$-algebra map to an $E_2$ refinement of
$R\operatorname{Hom}_A(P,P)$ can induce this characteristic
action. Such a map would preserve the Gerstenhaber bracket and
give
\[
 -\beta=\chi_P([E_u,\tau])
                  =[\chi_P(E_u),\chi_P(\tau)]=[0,\beta]=0.
\]
The ordinary derived characteristic action remains defined.
This obstruction concerns its specified values and that proposed
higher structure on its target.

\section{Adding constrained systems}

Relative interface composition uses one fixed algebra $A$.
Adding the coordinates of two independent constrained systems
asks instead for a coproduct compatible with their equations.
For the node this produces a cross term that is not a boundary
in the ordinary tensor of the two systems.

Consider the commutative dg algebra
\begin{equation}\label{cyrec:eq:constrained-node}
 K=k[u,m]\otimes\Lambda(\eta,\chi),\qquad
 |\eta|=|\chi|=-1,\quad d\eta=um,\quad d\chi=0.
\end{equation}
Its cohomology is $A\otimes\Lambda(\chi)$.
The proposed sums $u\mapsto u_1+u_2$ and
$m\mapsto m_1+m_2$ send the relation to
\[
 b_{12}=u_1m_2+u_2m_1
          \quad\bmod (u_1m_1,u_2m_2).
\]
This is nonzero, as setting $m_1=u_2=0$ shows.
Thus no dg algebra map $K\to K\otimes K$ can induce those
coordinate sums. Higher components of an $A_\infty$ map cannot
change the induced cohomology algebra relation either.

\subsection{A common refinement and a noninvertible comparison}

Adjoin an odd generator $\theta_{12}$ with
$d\theta_{12}=b_{12}$. The coupled algebra $T_2$ then receives
the map
\[
 u\mapsto u_1+u_2,\quad m\mapsto m_1+m_2,\quad
 \eta\mapsto\eta_1+\eta_2+\theta_{12},\quad
 \chi\mapsto\chi_1+\chi_2.
\]
The added generator makes it a chain map.
It does not yet make successive binary constructions equivalent.

For three leaves put $b_{ij}=u_im_j+u_jm_i$.
The left parenthesization has additional generators with
differentials $b_{12}$ and $b_{13}+b_{23}$.
The right parenthesization uses $b_{23}$ and $b_{12}+b_{13}$.
Their degree-zero cohomology algebras are the quotients of the
fixed six-coordinate polynomial ring by
\begin{equation}\label{cyrec:eq:node-support-ideals}
\begin{aligned}
 I_L&=(u_im_i\ (1\leq i\leq3),\ b_{12},\ b_{13}+b_{23}),\\
 I_R&=(u_im_i\ (1\leq i\leq3),\ b_{23},\ b_{12}+b_{13}).
\end{aligned}
\end{equation}

At $u=(1,-1,0)$ and $m=(0,0,1)$, every generator of $I_L$
vanishes but $b_{23}=-1$. Thus the ideals differ.
An equivalence over the fixed leaf-coordinate ring would have to
identify their kernels in degree-zero cohomology, which is impossible.

There is nevertheless a coherent common refinement. For $r$
leaves let
\[
 T_r=k[u_i,m_i]\otimes
       \Lambda(\eta_i,\chi_i,\theta_{ij}:i<j),
 \quad d\eta_i=u_im_i,\quad d\chi_i=0,\quad d\theta_{ij}=b_{ij}.
\]
For a partition into ordered blocks, the refinement map sends
\begin{equation}\label{cyrec:eq:node-refinements}
\begin{aligned}
 u_B&\mapsto\sum_{i\in B}u_i,&
 m_B&\mapsto\sum_{i\in B}m_i,\\
 \eta_B&\mapsto\sum_{i\in B}\eta_i+
                      \sum_{i<j\in B}\theta_{ij},&
 \chi_B&\mapsto\sum_{i\in B}\chi_i,\\
 \theta_{B,C}&\mapsto\sum_{i\in B,j\in C}\theta_{\{i,j\}}.
\end{aligned}
\end{equation}

Expansion proves compatibility with differentials. Every pair
occurs exactly once under successive refinement, so these maps
compose strictly. Both binary parenthesizations map to $T_3$.
Their maps to that common target are not thereby invertible
associators.

\subsection{Retaining the equation value gives a Hopf algebra}

Instead of imposing $um=0$, keep its value as a coordinate:
\[
 \widetilde K=k[u,m,c]\otimes\Lambda(\eta,\chi),
 \qquad d\eta=um-c,\quad d\chi=0.
\]
The coordinate $w=c-um$ makes the differential linear,
$d\eta=-w$.
Declare $u,m,w,\eta,\chi$ primitive in the free
graded-commutative algebra on these generators.

\begin{proposition}\label{cyrec:prop:equation-hopf}
This makes $\widetilde K$ a dg Hopf algebra, with
\begin{equation}\label{cyrec:eq:equation-coproduct}
 \Delta c=c\otimes1+1\otimes c+u\otimes m+m\otimes u,
 \qquad S(c)=-c+2um.
\end{equation}
Its counit kills every displayed generator. The quotient $c=0$
is not a Hopf quotient, and
$H^*(\widetilde K)=k[u,m]\otimes\Lambda(\chi)$.
\end{proposition}

\begin{proof}
In the variables $u,m,w,\eta,\chi$, the algebra is the free
graded-commutative Hopf algebra on a complex of primitives.
The linear differential commutes with all its Hopf maps.
Substitution $c=w+um$ gives both formulas in
\eqref{cyrec:eq:equation-coproduct}.
Modulo $c$ in both factors, the cross term remains nonzero,
so $(c)$ is not a coideal. Contracting the pair $(w,\eta)$
computes the cohomology.
\end{proof}

The enlargement provides a coherent tensor by retaining the
equation value, but changes the cohomology of the constrained
system. Relative interface composition, the common-refinement
maps, and this enlarged Hopf tensor therefore solve distinct
problems.

\begin{problem}\label{cyrec:prob:constrained-fusion}
Construct a category of constrained nodal field theories with a
specified tensor product, local vacuum and collision operations.
Determine whether it receives the associative interface
composition over $A$, or the all-pair refinement maps
\eqref{cyrec:eq:node-refinements}, through functors preserving
their derived coefficients. If an enlargement by equation values
is used, construct the comparison with the original constrained
cohomology and identify which operations it preserves.
\end{problem}
